% =====================================================================
%  capitulos/cap05.tex — Capítulo 5. Los constituyentes y la estructura
%  sintagmática
%  Autora: María Paula Bonorino
%  Fuente: Capitulo_05_normalizado.docx
% =====================================================================
\capitulo[Los constituyentes y la estructura sintagmática]%
  {Los constituyentes y la estructura sintagmática}{María Paula Bonorino}

\section{Introducción}

Este capítulo inicia el estudio de la sintaxis y para ello introduce una serie
de conceptos clave para reflexionar sobre las relaciones entre las palabras que
forman las oraciones, ya que, al igual que las unidades de la morfología, los
ítems léxicos no solo se ordenan de manera lineal, sino que también establecen
dependencias jerárquicas. En primer lugar, consideramos cómo se definen los
constituyentes sintácticos y presentamos las pruebas formales que permiten
identificarlos. En segundo lugar, nos centramos en las clases de palabras
léxicas, entendidas como aquellas que poseen significado descriptivo inherente,
y examinamos cómo se expanden en sintagmas, conformando así unidades intermedias
de análisis, que exhiben una organización interna. Finalmente, nos enfocamos en
las características de los diferentes tipos de sintagmas ---nominales (SSNN),
adjetivales (SSAA), adverbiales (SSADV) y preposicionales
(SSPP)\footnote{La duplicación de las siglas (SSNN, SSAA, etc.) indica
plural.}--- y consideramos sus principales propiedades sintácticas y
semánticas.

\section{El concepto de constituyente y la estructura jerárquica}\label{sec:05-constituyente}

Así como los morfemas se unen para formar palabras, también las palabras se
combinan en unidades mayores. Consideremos, por ejemplo, la siguiente oración:

\begin{exe}
  \ex\label{ex:05-ministro} El ministro insistió en su declaración.
\end{exe}

Si bien en \exref{ex:05-ministro} las palabras se vinculan linealmente, también
podemos reconocer ---al menos intuitivamente--- que la relación entre
\forma{ministro} y \forma{el} es más estrecha que la que existe entre
\forma{ministro} e \forma{insistió}, aunque estas dos últimas sean palabras
contiguas. La diferencia radica en que \forma{ministro} y \forma{el} están en el
mismo constituyente e \forma{insistió}, en cambio, se encuentra en otro. Pero
¿qué es un constituyente? Un \termino{constituyente} es un grupo de palabras que
funciona como una unidad «independientemente de si tiene una estructura interna
simple o compleja» \parencite[88]{camacho2018}. Así, en
\exref{ex:05-AB}, tenemos dos constituyentes mayores a los que llamaremos A y B:

\begin{exe}
  \ex\label{ex:05-AB} {}[El ministro]\etq{A} [insistió en su declaración]\etq{B}
\end{exe}

\begin{destacado}
El concepto de constituyente es relacional: «algo es un constituyente de una
determinada construcción si forma parte de esta\ldots{}»
\parencite[69]{ditullio2014}.
\end{destacado}

En línea con lo que plantea Di Tullio, un constituyente puede, a su vez, estar
integrado por constituyentes menores. El constituyente A está formado por el
determinante \forma{el} y el sustantivo \forma{ministro}.

\begin{exe}
  \ex {}[\etq{A} El ministro]
\end{exe}

El constituyente B, por su parte, plantea una estructura más compleja. Cabe
reconocer allí, en primer lugar, un constituyente formado por la combinación del
sustantivo \forma{declaración} con el posesivo \forma{su}. Este constituyente se
combina, a su vez, con la preposición \forma{en} para formar el constituyente
\forma{en su declaración}. Finalmente, este constituyente se combina con el verbo
\forma{insistió}, dando como resultado el constituyente B, \forma{insistió en su
declaración}.

\begin{exe}
  \ex {}[\etq{B} insistió [en [su declaración]]]
\end{exe}

\begin{destacado}
En síntesis, los constituyentes de una oración se relacionan entre sí de forma
jerárquica y no simplemente por contigüidad.
\end{destacado}

La estructura de constituyentes se puede representar mediante corchetes, como en
\exref{ex:05-corchetes}.

\begin{exe}
  \ex\label{ex:05-corchetes} {}[\etq{AB} [\etq{A} El ministro] [\etq{B} insistió
      [en [su declaración]]]]
\end{exe}

\begin{ejercicio}\label{ej:05-segmentar}
Segmente en constituyentes las siguientes oraciones, como se muestra en el
ejemplo resuelto~1.

\begin{exe}
  \ex Sus conclusiones asombraron a los científicos.
  \ex Nuestra tía traía regalos constantemente.
  \ex Muchos estudiantes prefieren la virtualidad.
  \ex Las medidas preocupan a la población.
  \ex Los jugadores ganaron la copa en 2022.
  \ex Aquellas profesoras incorporaron en sus clases herramientas tecnológicas.
\end{exe}

\begin{resuelto}
1. [[Sus conclusiones [sobre [el tema]]] [asombraron [a [la comunidad
científica]]]]\pendiente{La resolución no corresponde a la oración 1: «Sus
conclusiones asombraron a los científicos»}
\end{resuelto}
\end{ejercicio}

\subsection{Homonimia estructural}\label{sec:05-homonimia}

En algunos casos,\footnote{Este tema se presentó en el \capref{2},
§\ref{sec:02-ambiguedad}.} como se muestra en \exref{ex:05-salon}, la
segmentación en constituyentes puede ofrecer más de una alternativa
interpretativa: así, puede ser que \forma{el salón con flores} forme un solo
constituyente \exref{ex:05-salon-b} o que se trate de dos constituyentes
independientes, que pueden intercambiarse \exref{ex:05-salon-c}:

\begin{exe}
  \ex\label{ex:05-salon}
    \begin{xlist}
      \ex Redecoró el salón con flores.
      \ex\label{ex:05-salon-b} Redecoró [el salón [con flores]].
      \ex\label{ex:05-salon-c} Redecoró [el salón] [con flores]. / Redecoró [con
          flores] [el salón].
    \end{xlist}
\end{exe}

Las oraciones como \exref{ex:05-salon}, que admiten más de una segmentación en
constituyentes, son ejemplos de \termino{homonimia estructural}, el fenómeno
según el cual una oración presenta más de una descripción estructural (DE), cada
una asociada a una interpretación diferente de la oración. Así, en
\exref{ex:05-salon-b}, [\forma{el salón con flores}] constituye una unidad porque
hace referencia a \sig{un salón que ya tenía flores, por ejemplo, en el
empapelado}, y que fue redecorado, probablemente con otro diseño. En este caso,
[\forma{con flores}] es un constituyente menor que modifica a [\forma{el salón}],
por lo que, como constituyen una unidad, no admiten intercambiar sus posiciones.
En \exref{ex:05-salon-c}, en cambio, en la que interpretamos que \sig{puso flores
en el salón}, se puede considerar que [\forma{el salón}] y [\forma{con flores}]
son dos modificadores independientes del verbo \forma{redecoró} y que, por lo
tanto, sus posiciones son intercambiables.

\begin{destacado}
Dado que una misma oración puede recibir más de una descripción estructural,
según como se interpreten sus constituyentes, se considera que las secuencias
ambiguas establecen un argumento en favor de la teoría de que las construcciones
sintácticas poseen una estructura jerárquica.
\end{destacado}

\begin{ejercicio}\label{ej:05-homonimia}
Considere las siguientes oraciones y resuelva lo pedido, según se muestra en el
ejemplo resuelto~1.

\begin{consignas}
  \item Segmente en constituyentes.
  \item Proporcione las dos interpretaciones posibles de cada oración y explique
        a qué se deben.
  \item Luego de responder a los puntos a) y b), repase el
        ejercicio~\ref{ej:02-ambiguedad} del \capref{2} y considere qué efecto
        produce la doble posibilidad de segmentación que ofrecen estas
        oraciones.
\end{consignas}

\begin{exe}
  \ex Mi hermano trajo una foto de Egipto.
  \ex Llenó la bolsa de papel.
  \ex Encontraron el juguete roto.
\end{exe}

\begin{resuelto}
1. a) i) [[Mi hermano] [trajo [una foto] [de Egipto]]] / ii) [[Mi hermano]
[trajo [una foto de Egipto]]]. b) En i) se puede considerar que [una foto] y
[de Egipto] son dos modificadores independientes del verbo \forma{trajo}. En
ii), [de Egipto] modifica a [una foto], por lo que conforman un solo
constituyente.
\end{resuelto}
\end{ejercicio}

\subsection{Identificación de constituyentes}\label{sec:05-pruebas}

Para el reconocimiento de constituyentes se pueden aplicar pruebas que permiten
identificarlos. Las más habituales son las de sustitución, movimiento e
interrogación. Veamos en qué consiste cada una de ellas.

\subsubsection{Sustitución}\label{sec:05-sustitucion}

Toda unidad sintáctica compleja, como la secuencia subrayada en
\exref{ex:05-conicet-a}, puede concebirse como el desarrollo o expansión de otra
más simple por la que puede ser sustituida, como en \exref{ex:05-conicet-b} y
\exref{ex:05-conicet-c}, y toda unidad simple puede dar lugar a una expansión en
un constituyente más complejo, como se muestra en \exref{ex:05-conicet-d} y
\exref{ex:05-conicet-e}:

\begin{exe}
  \ex\label{ex:05-conicet}
    \begin{xlist}
      \ex\label{ex:05-conicet-a} \subrayado{Los investigadores del CONICET}
          analizaron con gran cuidado la nueva variante del virus.
      \ex\label{ex:05-conicet-b} \subrayado{Los investigadores} analizaron con
          gran cuidado la nueva variante del virus.
      \ex\label{ex:05-conicet-c} \subrayado{Ellos} analizaron con gran cuidado la
          nueva variante del virus.
      \ex\label{ex:05-conicet-d} \subrayado{Los investigadores más renombrados
          del CONICET} analizaron con gran cuidado la nueva variante del virus.
      \ex\label{ex:05-conicet-e} \subrayado{Los investigadores del CONICET
          (Consejo Nacional de Investigaciones Científicas y Técnicas)}
          analizaron con gran cuidado la nueva variante del virus.
    \end{xlist}
\end{exe}

\begin{ejercicio}\label{ej:05-expansion}
Proponga expansiones para los constituyentes subrayados, según se muestra en el
ejemplo~1.

\begin{exe}
  \ex \subrayado{Susana} firmó \subrayado{los papeles}.
  \ex \subrayado{El museo} recibió \subrayado{los ejemplares}.
  \ex \subrayado{Mis primas} \subrayado{llegaron}.
  \ex \subrayado{La temporada} trajo \subrayado{sorpresas}.
  \ex \subrayado{Los vehículos} circulaban \subrayado{lentamente}.
\end{exe}
\pendiente{Los subrayados no se conservan en el original; se reconstruyeron a
partir de las respuestas}

\begin{resuelto}[Ejemplo 1. Respuesta posible]
Susana, la directora del instituto, firmó los papeles necesarios para la
presentación del proyecto.
\end{resuelto}
\end{ejercicio}

\begin{ejercicio}\label{ej:05-sustitucion}
Sustituya en la oración los constituyentes subrayados. Proporcione para cada uno
de ellos una versión mediante sustitución por reducción y otra por
pronominalización, como se muestra para el ejemplo~1.

\begin{exe}
  \ex \subrayado{Los investigadores del CONICET} analizaron \subrayado{con gran
      cuidado} \subrayado{la nueva variante del virus}.
  \ex \subrayado{Varios ministros del gabinete} firmaron \subrayado{un acuerdo
      bilateral} \subrayado{con el embajador de la Unión Europea}.
\end{exe}

\begin{resuelto}
1. Reducción: Los investigadores analizaron con cuidado la variante.
Pronominalización: Ellos la analizaron así.
\end{resuelto}
\end{ejercicio}

\begin{ejercicio}\label{ej:05-encorchetar}
Considere las siguientes oraciones y resuelva las consignas a continuación,
según se muestra en el ejemplo resuelto~1.

\begin{consignas}
  \item Encorchete los constituyentes.
  \item Aplique la prueba de sustitución a todos los constituyentes
        encorchetados en a).
\end{consignas}

\begin{exe}
  \ex Sus tías enviaron los regalos por correo.
  \ex Los egresados buscaron con entusiasmo un salón de fiestas.
  \ex Nuestros candidatos presentaron sus propuestas a los contribuyentes.
  \ex Los fuertes temporales provocaron inundaciones en toda la ciudad.
\end{exe}

\begin{resuelto}
1. a) [[Sus tías] [enviaron [los regalos] [por correo]]]. b) Ellas enviaron los
regalos por correo. Sus tías los enviaron por correo. Sus tías enviaron los
regalos así. Ellas los enviaron así.
\end{resuelto}
\end{ejercicio}

\subsubsection{Interrogación}\label{sec:05-interrogacion}

Solo los constituyentes pueden funcionar como respuesta a preguntas formuladas
con una palabra interrogativa como \forma{qué, quién, cómo, cuándo}, etc. Estas
preguntas parciales\footnote{Las preguntas pueden ser totales o parciales:
mientras las parciales se responden mediante un constituyente de la oración,
como se muestra en §\ref{sec:05-interrogacion}, las preguntas totales son las
que interrogan por todo el enunciado y se responden afirmativa o negativamente:
\forma{¿Querés venir hoy al cine conmigo? Sí / No (quiero)}.} tienen que
responderse mediante algún constituyente específico de la oración, lo que
permite probar su condición de tal. Así, por ejemplo, respecto de
\exref{ex:05-conicet-a}, podríamos preguntar:

\begin{exe}
  \ex\label{ex:05-preguntas}
    \begin{xlist}
      \ex \forma{¿Quiénes} analizaron con gran cuidado la nueva variante del
          virus? Los investigadores del CONICET.
      \ex \forma{¿Qué} hicieron los investigadores del CONICET? Analizaron con
          gran cuidado la nueva variante del virus.
      \ex \forma{¿Cómo} analizaron los investigadores del CONICET la nueva
          variante del virus? Con gran cuidado.
      \ex \forma{¿Qué} analizaron los investigadores del CONICET con gran
          cuidado? La nueva variante del virus.
    \end{xlist}
\end{exe}

Como se advierte en los ejemplos anteriores, cada pregunta parcial recibe como
respuesta un constituyente de la oración.

\begin{ejercicio}\label{ej:05-preguntas}
¿Qué preguntas podríamos hacer para obtener cada uno de los constituyentes
encorchetados en la descripción estructural de las siguientes oraciones?

\begin{exe}
  \ex {}[El cartero] [trajo [a mi casa] [la nueva tarjeta de crédito]]
  \ex {}[Mis primos] [llegaron [ayer por la tarde] [a la ciudad] [para
      visitarnos]]
  \ex {}[La nueva profesora] es [muy exigente] [con los exámenes]
\end{exe}

\begin{resuelto}
1. ¿Quién trajo a mi casa la nueva tarjeta de crédito? / ¿Qué hizo el cartero? /
¿Adónde trajo el cartero la nueva tarjeta de crédito? / ¿Qué trajo a mi casa el
cartero?
\end{resuelto}
\end{ejercicio}

\begin{ejercicio}\label{ej:05-bioy}
Recurra a la prueba de la interrogación para verificar que las partes subrayadas
son constituyentes.

% Fragmento literario protegido: se incorpora a mano desde el original.
\textoaincorporar{el comienzo de «La salvación», de A.~Bioy Casares
(\textit{La invención y la trama}, Tusquets, 1991), con los constituyentes
subrayados: \textit{el escultor; con el tirano; por los jardines del palacio;
más allá del laberinto para los extranjeros ilustres, en el extremo de la
alameda de los filósofos decapitados; su última obra}}

\begin{resuelto}
¿Quién paseaba con el tirano por los jardines del palacio?
\end{resuelto}
\end{ejercicio}

\subsubsection{Movimiento}\label{sec:05-movimiento}

Los constituyentes se desplazan en bloque, por lo que solo un constituyente
completo puede moverse a otra posición dentro de la oración. En
\exref{ex:05-conicet-a}, que repetimos aquí como \exref{ex:05-mov-a}, los dos
constituyentes mayores pueden intercambiar sus posiciones, dando lugar a
\exref{ex:05-mov-b}:

\begin{exe}
  \ex\label{ex:05-mov}
    \begin{xlist}
      \ex\label{ex:05-mov-a} {}[[Los investigadores del CONICET] [analizaron con
          gran cuidado la nueva variante del virus]].
      \ex\label{ex:05-mov-b} {}[[Analizaron con gran cuidado la nueva variante
          del virus] [los investigadores del CONICET]].
    \end{xlist}
\end{exe}

Al realizar este movimiento, la oración sigue siendo gramatical. Asimismo,
también se puede hacer la misma operación con otros constituyentes presentes en
la oración:

\begin{exe}
  \ex {}[Con gran cuidado] [los investigadores del CONICET] [analizaron [la nueva
      variante del virus]].
\end{exe}

Sin embargo, si se delimita erróneamente un constituyente, el movimiento puede
dar lugar a una oración agramatical:

\begin{exe}
  \ex[*]{\label{ex:05-mov-agram} La nueva variante [los investigadores del
      CONICET] [analizaron [del virus] [con gran cuidado]].}
\end{exe}

En \exref{ex:05-mov-agram}, el constituyente \forma{la nueva variante del virus}
se ha fragmentado, lo que da como resultado una oración agramatical en español.

\begin{ejercicio}\label{ej:05-pares}
Compare los siguientes pares de oraciones y luego, según se muestra en el
ejemplo resuelto~1:

\begin{consignas}
  \item indique mediante encorchetamiento si las expresiones en cursiva
        conforman o no un único constituyente;
  \item fundamente la respuesta a partir de las pruebas de reconocimiento
        estudiadas.
\end{consignas}

\begin{exe}
  \ex
    \begin{xlist}
      \ex Le regalaron \textit{un pasaje a Daniel}.
      \ex Le regalaron \textit{un pasaje a México}.
    \end{xlist}
  \ex
    \begin{xlist}
      \ex Eliana encontró \textit{el auto de repente}.
      \ex Eliana encontró \textit{el auto de su primo}.
    \end{xlist}
  \ex
    \begin{xlist}
      \ex El hombre tocó \textit{al perro con la mano}.
      \ex El hombre tocó \textit{al perro con sarna}.
    \end{xlist}
  \ex
    \begin{xlist}
      \ex Juan abrió \textit{la puerta de una patada}.
      \ex Juan abrió \textit{la puerta de la habitación}.
    \end{xlist}
\end{exe}

\begin{resuelto}
1. a) i. Le regalaron [un pasaje] [a Daniel] / ii. Le regalaron [un pasaje a
México]. b) En i., \forma{un pasaje} y \forma{a Daniel} son dos constituyentes
que admiten las tres pruebas propuestas: se mueven separadamente (\forma{A Daniel
le regalaron un pasaje. Un pasaje le regalaron a Daniel}); responden a distintas
preguntas (\forma{¿Qué le regalaron a Daniel? ¿A quién le regalaron un pasaje?});
se pronominalizan independientemente (\forma{Le regalaron un pasaje. Se lo
regalaron a Daniel}).\footnote{V. la nota~\ref{fn:05-se}.} En
ii., \forma{un pasaje a México} es un solo constituyente porque se mueve
conjuntamente (\forma{Un pasaje a México le regalaron}); responde a una única
pregunta (\forma{¿Qué le regalaron?}); admite pronominalización conjunta
(\forma{Se lo regalaron}).
\end{resuelto}
\end{ejercicio}

\section{Los sintagmas}\label{sec:05-sintagmas}

Como vimos cuando aplicamos la prueba de sustitución, un constituyente puede
reducirse a un único elemento central: el núcleo. Por ejemplo, en
\exref{ex:05-pedro}, ambos constituyentes de la oración están formados por un
solo elemento:

\begin{exe}
  \ex\label{ex:05-pedro} {}[Pedro]\etq{A} [corre]\etq{B}
\end{exe}

Sin embargo, por lo general, los núcleos se expanden vinculándose con otras
palabras y conformando unidades mayores denominadas \termino{sintagmas}. Desde
esta perspectiva, \textcite[138]{bosque2009} definen los sintagmas como
«expansiones de un elemento central o nuclear, en torno al que se articulan
modificadores y complementos».

\begin{destacado}
Según vemos, entonces, todos los sintagmas tienen núcleo, es decir, son
endocéntricos y, en todos ellos, es el elemento nuclear el que da nombre al
sintagma.
\end{destacado}

\subsection{Las clases de palabras y su proyección sintagmática}\label{sec:05-clases}

Las palabras de la lengua se agrupan en dos clases mayores: las léxicas, por un
lado, y las gramaticales o funcionales, por el otro. Palabras como \forma{libro,
alto, leer, fácilmente} o \forma{para} se consideran léxicas porque su contenido
puede asociarse con ideas o conceptos \parencite[52]{rodriguezramalle2005}. Las
\termino{palabras léxicas} o de contenido descriptivo inherente incluyen el
sustantivo (N), el adjetivo (A), el verbo (V), el adverbio (Adv) y la
preposición (P).\footnote{Respecto de la preposición como palabra léxica, v. más
adelante §\ref{sec:05-sp}.} Pero, además de estas clases, también integran la
lengua las que Bosque (1991)\pendiente{En la bibliografía figura Bosque (1989)}
denominó «palabras engranaje», como \forma{y, mi, el} o \forma{que}, que no
tienen significado en sí mismas, sino que se definen fundamentalmente por la
función que cumplen en la oración. Son las llamadas \termino{palabras
funcionales}, constituidas por los determinativos (artículos, demostrativos,
posesivos, indefinidos, etc.) y las conjunciones (coordinantes y
subordinantes).

Aunque más adelante, en el \capref{6}, se verá que también las palabras y
categorías funcionales proyectan sintagmas en los que funcionan como núcleos,
aquí nos vamos a concentrar en los sintagmas formados por las categorías
léxicas. Todas las palabras léxicas se expanden o proyectan en sintagmas de los
que son núcleo, de modo que tenemos sintagmas nominales (SSNN), sintagmas
adjetivales (SSAA), sintagmas adverbiales (SSADV), sintagmas preposicionales
(SSPP) y sintagmas verbales (SSVV). En este capítulo nos ocupamos solo de las
cuatro primeras clases de sintagmas, ya que, dada la complejidad del SV, sus
particularidades se tratarán con más detalle en los capítulos
siguientes.\footnote{En el \capref{6} se trata la oración y su estructura; en el
\capref{7}, la estructura argumental y sintáctica, y en el \capref{8}, la
estructura eventiva.}

\begin{ejercicio}\label{ej:05-nucleo}
En los siguientes sintagmas, según se muestra en el ejemplo resuelto~1:

\begin{consignas}
  \item identifique el núcleo (tenga en cuenta cuál es la palabra central a la
        que las demás modifican o complementan);
  \item indique la categoría del núcleo de cada sintagma y señale qué clase de
        sintagma proyecta.
\end{consignas}

\begin{exe}
  \ex Ayer por la tarde.
  \ex Bueno con la gente.
  \ex Cerca de la esquina.
  \ex Directamente proporcional.
  \ex Para mayores de edad.
  \ex Profunda inquietud.
  \ex Armario antiguo.
  \ex Sin sentimientos negativos.
\end{exe}

\begin{resuelto}
1. a) El núcleo es \forma{ayer}. b) \forma{Ayer} es un adverbio y proyecta un
SAdv.
\end{resuelto}
\end{ejercicio}

\subsection{La estructura del sintagma}\label{sec:05-xbarra}

De acuerdo con el modelo de principios y parámetros propuesto por
\textcite{chomsky1981} (v. \capref{1}, §\ref{sec:01-estructura}), la teoría
X\prima{} o teoría de la X con barra, que se encarga de la buena formación de
los sintagmas, plantea que todos los sintagmas de todas las lenguas poseen la
misma estructura. No solo son siempre endocéntricos, es decir, poseen un núcleo
que da la categoría a todo el sintagma, sino que también adoptan la misma
configuración estructural. Veamos un ejemplo:

\begin{exe}
  \ex\label{ex:05-presidente} {}[El presidente de la república] inauguró las
      sesiones legislativas.
\end{exe}

Si consideramos el SN \forma{el presidente de la república} y tenemos en cuenta
cómo se conforman los constituyentes, podemos decir que, dado que \forma{ser
presidente} exige \forma{presidir algo}, el núcleo \forma{presidente} primero se
une a \forma{de la república}, que, al ser un constituyente exigido por el
núcleo, funciona como su complemento (Comp). El tercer elemento que integra el
SN es el artículo \forma{el}, también requerido para la buena formación de la
oración, ya que \juicio{*}{\forma{Presidente de la república inauguró las
sesiones legislativas}}\footnote{Dicha oración sería aceptable, pero en un
registro muy particular, como son los titulares de diarios.} no sería gramatical
en español. Por tanto, ambos, núcleo y complemento, deben ir precedidos del
artículo \forma{el}, que los determina y cumple la función de especificador
(Esp) dentro del sintagma. Podemos representar la estructura sintagmática que
acabamos de describir mediante encorchetamiento \exref{ex:05-presidente-a} o a
través de un diagrama arbóreo \exref{ex:05-presidente-b}:

\begin{exe}
  \ex\label{ex:05-presidente-de}
    \begin{xlist}
      \ex\label{ex:05-presidente-a} {}[\etq{SN} el\etq{D/Esp} [\etq{N\prima}
          presidente\etq{N/n} [\etq{SP/Comp} de la república]]]
      \ex\label{ex:05-presidente-b}
          \begin{forest} gramatica
            [SN
              [D\\\textit{el}]
              [N\prima
                [N\\\textit{presidente}]
                [SP [\textit{de la república}, roof]]]]
          \end{forest}
    \end{xlist}
\end{exe}

La estructura propuesta no es plana, es decir, no todos los elementos están en
el mismo nivel. Según se puede ver en \exref{ex:05-presidente-b}, la estructura
del sintagma es jerárquica e implica niveles que describen la relación entre los
constituyentes. En primer lugar, en el grado X se ubica el núcleo, que es el
elemento que no puede faltar, define la categoría del sintagma y selecciona los
complementos y modificadores. En el segundo nivel, X\prima, el núcleo se une a
su complemento, que es el elemento con el que tiene una relación más estrecha y,
por último, en el tercer nivel, el especificador es un constituyente que, según
su clase, determina ---\forma{el / este / mi presidente}--- o cuantifica
---\forma{algún / varios / todos los presidentes}--- conjuntamente al núcleo y a
su complemento. El sintagma es, pues, la máxima proyección del núcleo y
constituye el nivel X\prima\prima{} o SX (sintagma de X):

\begin{exe}
  \ex\label{ex:05-esquema}
      \begin{forest} gramatica
        [SX (X\prima\prima)
          [Esp]
          [X\prima
            [X\\(núcleo)]
            [Comp]]]
      \end{forest}
\end{exe}

\subsubsection{Complementos, especificadores y adjuntos}\label{sec:05-comp-esp-adj}

Dado que todas las categorías presentan la misma estructura, este esquema puede
aplicarse a cualquiera de ellas, si se sustituye la variable X por una etiqueta
categorial, como podrían ser N, V, A, P o Adv.\footnote{En este capítulo nos
centramos en los núcleos formados por palabras léxicas ---N, A, Adv y P---,
excepto el V, que se verá en el \capref{6}. A partir del \capref{7} nos iremos
adentrando en los sintagmas proyectados por palabras y categorías
funcionales.}

Es necesario remarcar que hay categorías transitivas, como algunos verbos y como
todas las preposiciones, que siempre exigen un complemento, y categorías
intransitivas, como los adverbios y algunos verbos. En las categorías
intransitivas, la posición de complemento queda vacía, como se muestra en
\exref{ex:05-sonrio}, en el que aparece el verbo intransitivo
\forma{sonreír}.\footnote{Las distintas clases de verbos se presentan en el
\capref{7}.} Asimismo, la posición de especificador también puede quedar vacía,
como se muestra en \exref{ex:05-irene}, donde el núcleo del SN sujeto es un
nombre propio, que no requiere especificador.\footnote{V. más adelante
§\ref{sec:05-sn}.}

\begin{exe}
  \ex\label{ex:05-sonrio} El muchacho [\etq{SV} sonrió\etq{V/n}].
  \ex\label{ex:05-irene} {}[\etq{SN} Irene\etq{N/n}] vendrá a visitarnos.
\end{exe}

Dentro de la estructura sintagmática, los \termino{complementos} son
constituyentes seleccionados o exigidos por el núcleo en función de sus
requerimientos léxicos y sintácticos. Así, por ejemplo, un verbo transitivo,
como \forma{comer}, \forma{escribir} o \forma{leer}, exige un SN vinculado con
su significado que lo complemente: \forma{comió torta, escribió una carta, leyó
un libro}. Desde el punto de vista estructural, también las preposiciones exigen
un complemento, ya que no pueden funcionar solas: \forma{de madera, por tu
culpa, sin aliento}.

Aunque los complementos son exigidos por el núcleo, esto no quiere decir que
siempre tengan que estar explicitados en la sintaxis, ya que pueden quedar
implícitos; por ejemplo, en \forma{Hoy comimos temprano} se sobreentiende que se
ingirió \sig{algo comestible}, que no se desea especificar.

Los \termino{especificadores} modifican conjuntamente al núcleo junto con sus
complementos y modificadores. La posición del especificador suele estar ocupada
por palabras funcionales o gramaticales, como los determinantes [D] (\forma{el /
este / su libro}) o los cuantificadores [Q] (\forma{varias / ambas / algunas
revistas}, \forma{bastante / muy desgastado}).\footnote{Esta información se
amplía más adelante, en §\ref{sec:05-sn}.}

Además, los sintagmas admiten constituyentes más externos, no exigidos por el
núcleo, aunque compatibles con él, como los \termino{adjuntos}, que pueden
modificar a cualquier tipo de unidad \parencite[49]{rodriguezramalle2005}. Con
la incorporación del adjunto, la estructura sintagmática queda como se muestra
en \exref{ex:05-adjunto}:

\begin{exe}
  \ex\label{ex:05-adjunto}
      \begin{forest} gramatica
        [SX
          [Esp]
          [X\prima
            [X\prima
              [X]
              [Comp]]
            [Adjunto]]]
      \end{forest}
\end{exe}

Los adjuntos son opcionales, es decir, su presencia no es obligatoria dentro del
sintagma, como es el caso de los circunstanciales dentro del SV (\forma{conversaron
en la oficina / llovió torrencialmente / se encendió con alcohol}) o de los
modificadores dentro del SN (\forma{camisa roja / mantel redondo / aulas pequeñas
/ luces de colores / fragancias de París}).

En \exref{ex:05-dibujo} el verbo \forma{hacer} está acompañado de dos
constituyentes, [\forma{un dibujo}] y [\forma{en su cuaderno}], pero solo
\forma{un dibujo}, que explicita \sig{lo hecho en el cuaderno}, es exigido por el
núcleo verbal, ya que su ausencia (\juicio{*}{\forma{Hizo en su cuaderno}}) torna
la secuencia agramatical. En cambio, si se omite \forma{en su cuaderno}, la
oración «pierde» información, pero se mantiene bien formada, lo que indica que
ese constituyente es un adjunto.

\begin{exe}
  \ex\label{ex:05-dibujo}
    \begin{xlist}
      \ex Hizo [un dibujo] [en su cuaderno].
      \ex Hizo [un dibujo].
      \ex[*]{Hizo [en su cuaderno].}
    \end{xlist}
\end{exe}

\begin{ejercicio}\label{ej:05-de}
Realice la descripción estructural (DE) de los siguientes sintagmas mediante
encorchetamiento, según se muestra en el ejemplo resuelto~1.

\begin{exe}
  \ex Muy afecto al baile.
  \ex Varios colaboradores expertos.
  \ex Arriba de la heladera.
  \ex Sin comentarios.
\end{exe}

\begin{resuelto}
1. [\etq{SA} muy\etq{Q/Esp} afecto\etq{A/n} [\etq{SP/Comp} a\etq{P/n}
[\etq{SN/Comp} el\etq{D/Esp} baile\etq{N/n}]]]
\end{resuelto}
\end{ejercicio}

\begin{ejercicio}\label{ej:05-comp-adj}
Considere los siguientes sintagmas y luego, según se muestra en el ejemplo
resuelto~1:

\begin{consignas}
  \item identifique el núcleo e indique el tipo de sintagma que proyecta;
  \item distinga si los constituyentes que acompañan a los núcleos son
        complementos o adjuntos. Justifique su respuesta de acuerdo con la clase
        de palabra a la que pertenece cada núcleo y sus requerimientos para que
        la secuencia quede bien formada.
\end{consignas}

\begin{exe}
  \ex Orgullosa de sus logros.
  \ex La feria del libro.
  \ex Similar a los demás casos.
  \ex Por la izquierda.
  \ex Bajo la mesada.
  \ex Antes de las diez.
\end{exe}

\begin{resuelto}
1. a) El núcleo del sintagma es \forma{orgullosa} y proyecta un SA. b) \forma{De
sus logros} es un SP exigido por \forma{orgullosa} y funciona como Comp.
\end{resuelto}
\end{ejercicio}

\begin{ejercicio}\label{ej:05-titulos}
Considere los títulos de las obras que se presentan a continuación y resuelva
las consignas siguientes, según se muestra en los ejemplos resueltos:

\begin{consignas}
  \item identifique el núcleo y clasifique cada sintagma;
  \item determine la función que cumplen los constituyentes que acompañan al
        núcleo (especificador, complemento o adjunto). Justifique su respuesta
        considerando lo visto anteriormente.
\end{consignas}

\begin{exe}
  \ex \textit{Caos} (Magalí Tajes)
  \ex \textit{Un mundo feliz} (Aldous Huxley)
  \ex \textit{El jardín de los cerezos} (Antón Chéjov)
  \ex \textit{Cien años de soledad} (Gabriel García Márquez)
  \ex \textit{Ágilmente} (Estanislao Bachrach)
  \ex \textit{Ese verano a oscuras} (Mariana Enriquez)
  \ex \textit{Terrenal} (Mauricio Kartun)
  \ex \textit{Desde la ventana} (Carmen Martín Gaite)
\end{exe}

\begin{resuelto}[Ejemplos resueltos]
1. a) Es un SN porque el núcleo es el sustantivo \forma{caos}. b) El núcleo no
presenta complementos ni modificadores. 2. a) Es un SN porque el núcleo es el
sustantivo \forma{mundo}. b) \forma{Feliz} es un modificador adjunto porque no
está seleccionado por el núcleo.
\end{resuelto}
\end{ejercicio}

\begin{ejercicio}\label{ej:05-vf}
Integre los conocimientos incorporados hasta ahora y decida si las siguientes
afirmaciones son verdaderas o falsas, según se muestra en el ejemplo
resuelto~1. Fundamente sus respuestas.

\begin{exe}
  \ex Las palabras léxicas ---N, A, Adv, V y P--- se expanden o proyectan en
      SSXX.
  \ex El núcleo de un sintagma es siempre una palabra gramatical.
  \ex Todos los constituyentes de un sintagma se vinculan de la misma forma con
      el núcleo.
  \ex Los complementos son requeridos por el núcleo y se diferencian de los
      adjuntos, que son opcionales.
\end{exe}

\begin{resuelto}
1. V(erdadero). Las palabras léxicas son los núcleos de los sintagmas que
proyectan.
\end{resuelto}
\end{ejercicio}

\section{Las clases de sintagmas}\label{sec:05-clasesintagmas}

En los apartados que siguen, analizamos con mayor detenimiento las
características de los distintos tipos de sintagmas.

\subsection{El sintagma nominal (SN) y su estructura interna}\label{sec:05-sn}

Como ya hemos visto, el SN es aquel que tiene por núcleo a un sustantivo. En
cuanto a su estructura interna, como sostiene \textcite[153]{ditullio2014},
«resulta muy compleja debido a las múltiples posibilidades de realización que
presentan sus dependencias y a la estrecha conexión entre estas y el tipo de
sustantivo que funciona como núcleo».

De acuerdo con la teoría X\prima, que presentamos en §\ref{sec:05-xbarra}, el SN
posee la siguiente estructura:

\begin{exe}
  \ex\label{ex:05-esquema-sn}
      \begin{forest} gramatica
        [SN
          [Esp\\D / Q]
          [N\prima
            [N\prima
              [N]
              [Comp\\SP]]
            [Mod\\SA / SP / SN]]]
      \end{forest}
\end{exe}

\paragraph{El núcleo.} Cualquier clase semántica de sustantivo puede ser núcleo
de un SN, tanto común \exref{ex:05-armario} como propio \exref{ex:05-martin},
colectivo \exref{ex:05-bandada}, individual concreto \exref{ex:05-armario} o
abstracto \exref{ex:05-ilusion}.\pendiente{El original remite a (23) para los
abstractos}

\begin{exe}
  \ex\label{ex:05-armario}
    \begin{xlist}
      \ex El \textit{armario} del fondo.
      \ex Un \textit{edificio} muy alto del centro de la ciudad.
    \end{xlist}
  \ex\label{ex:05-martin}
    \begin{xlist}
      \ex \textit{Martín}, mi primo de Catamarca.
      \ex \textit{Luisa}, la profesora de música.
    \end{xlist}
  \ex\label{ex:05-bandada}
    \begin{xlist}
      \ex Una \textit{bandada} de pájaros.
      \ex\label{ex:05-arboleda} Esa lejana \textit{arboleda} de mi infancia.
    \end{xlist}
  \ex\label{ex:05-ilusion}
    \begin{xlist}
      \ex Una \textit{ilusión} pasajera como el verano.
      \ex El \textit{consentimiento} mutuo de la pareja.
    \end{xlist}
\end{exe}

En el SN, el núcleo determina los rasgos de género y número de los modificadores
y especificadores. Así, en \exref{ex:05-arboleda}, tanto \forma{lejana} como
\forma{esa} están en femenino singular para concordar con el núcleo
\forma{arboleda}.

\paragraph{Los complementos.} Si bien la mayor parte de los nombres, como
\forma{mesa} o \forma{cuaderno}, no tienen como requisito semántico la exigencia
de un complemento que complete su significación, hay algunos sustantivos
específicos que requieren un SP como complemento, como los siguientes:

\begin{itemize}
  \item Los \termino{nombres eventivos}, que son aquellos que por su sentido
        implican la ocurrencia de un suceso, que se desarrolla en un tiempo y
        lugar. Muchos nombres eventivos son deverbales, es decir, provienen de
        verbos (\forma{construcción, agasajo, reunión}), pero otros no
        (\forma{fiesta, batalla, partido}). Se reconocen porque pueden ser
        sujeto de \forma{tener lugar} \exref{ex:05-lugar} y complemento de la
        preposición \forma{durante} \exref{ex:05-durante}:
\end{itemize}

\begin{exe}
  \ex\label{ex:05-lugar} {}[La \textit{formulación} de una nueva teoría] tuvo
      lugar en la Universidad de Buenos Aires.
  \ex\label{ex:05-durante} El suceso ocurrió durante [la \textit{guerra} de los
      seis días].
\end{exe}

\begin{itemize}
  \item Los \termino{nombres relacionales}, que son los que establecen un vínculo
        semántico con lo designado por el complemento, como ser \forma{padre de
        alguien, alumno de algún profesor} o \forma{especialista en alguna
        disciplina}. Son ejemplos de SSNN con nombres relacionales los
        siguientes:
\end{itemize}

\begin{exe}
  \ex {}[El \textit{hijo} de Susana] vino a visitarnos.
  \ex {}[El \textit{autor} de esa novela] vive en Francia.
\end{exe}

\begin{itemize}
  \item Los \termino{nombres de representación}, que son los que designan una
        imagen o descripción de lo manifestado por el complemento. Son ejemplos
        de SSNN de representación:
\end{itemize}

\begin{exe}
  \ex Esa es [la \textit{foto} de la abuela] cuando era joven.
  \ex {}[La \textit{imagen} de un héroe de Malvinas] me impactó.
\end{exe}

En todos los casos, los complementos seleccionados también pueden quedar
implícitos:

\begin{exe}
  \ex Puse en un marco plateado [el \textit{retrato} (de la
      abuela)\footnote{Convencionalmente, los elementos opcionales se colocan
      entre paréntesis.}].
\end{exe}

\paragraph{Los modificadores adjuntos.} Son constituyentes opcionales que
aportan características del nombre y permiten restringirlo a un subconjunto
determinado de la clase general de entidades que abarca. Los modificadores
típicos del nombre son los adjetivos, tanto los calificativos, que indican
propiedades, \exref{ex:05-hojas} y \exref{ex:05-fornidos}, como los
relacionales, que vinculan al nombre con un ámbito determinado. Dado que los
relacionales son derivados de sustantivos, la relación también puede expresarse
mediante un SP, \exref{ex:05-tambera} y \exref{ex:05-lunar}
(v.~§\ref{sec:05-sa}):

\begin{exe}
  \ex\label{ex:05-hojas} Las hojas \textit{rayadas} de mi cuaderno \textit{azul}.
  \ex\label{ex:05-fornidos} Los \textit{fornidos} trabajadores de las minas.
  \ex\label{ex:05-tambera} Una rica zona \textit{tambera / de tambos}.
  \ex\label{ex:05-lunar} El inmenso cráter de la superficie \textit{lunar / de la
      luna}.
\end{exe}

Otro tipo de modificador del nombre es la \termino{aposición}, formada por un
sustantivo que pertenece a la misma clase semántica del núcleo cuando está en
relación de adyacencia directa con él, como vemos en \exref{ex:05-gregorio} y
\exref{ex:05-gomez}.

\begin{exe}
  \ex\label{ex:05-gregorio} {}[Gregorio, \textit{mi gato},] está todo el día
      tirado en la cocina.
  \ex\label{ex:05-gomez} {}[El comisario \textit{Gómez}] ya está retirado del
      servicio activo.
\end{exe}

\begin{recuadro}[Para ampliar: modificadores restrictivos y explicativos]
Los modificadores adjetivos y las aposiciones pueden ser restrictivos o
explicativos. Los restrictivos, también llamados especificativos, acotan la
referencia del nombre \exref{ex:05-andino}; los explicativos, en cambio, agregan
información complementaria, por lo que van marcados prosódicamente por una
figura tonal propia, lo que se refleja en la escritura porque se colocan entre
comas y se pueden suprimir \exref{ex:05-emilce}.

\begin{exe}
  \ex\label{ex:05-andino}
    \begin{xlist}
      \ex {}[El paisaje \textit{andino}] me pareció espectacular.
      \ex {}[Los hermanos \textit{Grimm}] le agregaron un final \textit{feliz} a
          la historia. (Adaptado de A.~Litvin, \textit{El libro de los
          villanos}, V\&R, 2015)
    \end{xlist}
  \ex\label{ex:05-emilce}
    \begin{xlist}
      \ex {}[Emilce, \textit{agitada},] señaló con su dedito de tres años la
          puerta abierta de su cuarto. (M.~Di Marco, \textit{El fantasma del
          Reich}, Sudamericana, 1995)
      \ex {}[Jacobo, \textit{el niño tonto},] solía subirse a la azotea y espiar
          la vida de los vecinos. (E.~Anderson Imbert, \textit{El gato de
          Cheshire}, Losada, 1965)
    \end{xlist}
\end{exe}

También los SSPP pueden modificar al nombre y son, en general, modificadores
restrictivos:

\begin{exe}
  \ex
    \begin{xlist}
      \ex Ahora trabaja en [una escuela \textit{con ruralidad}].
      \ex {}[La entrevista \textit{de mañana}] va a ser fundamental para el
          proyecto.
    \end{xlist}
\end{exe}
\end{recuadro}

\paragraph{Los especificadores.} Los sustantivos comunes suelen ir precedidos
por un determinativo en la posición de especificador. Los determinativos son
palabras funcionales que se definen por su relación con otra u otras palabras de
las que delimitan su alcance o a las que cuantifican al indicar «cuáles o
cuántos de los elementos incluidos en la clase denotada [\ldots] se deben
considerar» \parencite[93]{rodriguezramalle2005}. La clase de los determinativos
está formada por los determinantes [D] y los cuantificadores [Q], que son
palabras que carecen de contenido descriptivo específico y no establecen una
relación semántica con el sustantivo al que preceden, por lo que pueden
anteponerse a cualquier tipo de sustantivo, del que anticipan el género y el
número:

\begin{exe}
  \ex {}[\textit{El / nuestro / este / algún} pueblo] no debe estar lejos de
      aquí.
  \ex {}[\textit{Mis / estas / varias} amigas] vendrán a visitarme para mi
      cumpleaños.
\end{exe}

La subclase de los determinantes está formada por los artículos (\forma{el, los,
la, las, lo}), los demostrativos (\forma{este, esos, aquella}) y los posesivos
(\forma{mi, nuestros, tus}). Los cuantificadores, por su parte, están conformados
por los numerales cardinales (\forma{cinco, ocho, cien}), los numerales ordinales
(\forma{primero, segundo, decimoquinto}) y los indefinidos (\forma{mucho/s,
poco/s, varios/as, ambos/as, sendos/as, cualquiera, todo/s, bastante,
demasiado}).

A diferencia de los sustantivos comunes, que requieren un determinativo para
funcionar oracionalmente, los propios, por lo general, se emplean solos como
núcleo del SN. Los nombres propios no necesitan especificador porque ya están
determinados; así, en español no es necesario especificar \forma{la Irene, el
Pedro}. Sin embargo, en varios dialectos los hablantes colocan un especificador
delante de los nombres propios para indicar familiaridad o que se trata de
alguien conocido:

\begin{exe}
  \ex
    \begin{xlist}
      \ex La Frida\ldots{} Qué linda estás. (R.~Cossa, \textit{Gris de
          ausencia}, 1981)
      \ex ¿Y vos sabés cómo sufren los yanquis por eso? ¡Uf! Todos los días se
          reúnen dos horas a llorar porque la Yoly de Lanús los tiene acá.
          (N.~Fernández Tiscornia, \textit{Made in Lanús}, Cántaro, 1986)
    \end{xlist}
\end{exe}

\begin{ejercicio}\label{ej:05-buenaformacion}
Considere la buena o mala formación de los SSNN siguientes. Tenga en cuenta la
clase semántica de los sustantivos que funcionan como núcleo y sus posibilidades
y exigencias combinatorias, según se muestra en el ejemplo resuelto~1.

\begin{exe}
  \ex Trajo un regalo para novio.
  \ex El Mario.
  \ex Los ingredientes de la receta.
  \ex Compraron arena.
  \ex Vive con alegría.
\end{exe}

\begin{resuelto}
1. Es agramatical porque \forma{novio} es un sustantivo común y, por lo tanto,
debe ir precedido por un determinativo.
\end{resuelto}
\end{ejercicio}

\begin{ejercicio}\label{ej:05-ratones}
Lea detenidamente el siguiente texto, considere los SSNN subrayados y señale sus
constituyentes mediante corchetes, según se muestra en el ejemplo resuelto.

% Texto de divulgación protegido: se incorpora a mano desde el original.
\textoaincorporar{«A los ratones no les gusta el queso», de A.~Litvin
(\textit{El libro de las mentiras}, V\&R, 2015), con los SSNN subrayados:
\textit{los ratones; el queso; una compañía de control de plagas; los alimentos
con alto contenido de azúcares}}

\begin{resuelto}
{}[\etq{SN} los\etq{D/Esp} ratones\etq{N/n}]
\end{resuelto}
\end{ejercicio}

\begin{ejercicio}\label{ej:05-titulares}
Lea con atención los siguientes titulares de diarios con sus respectivas
volantas\footnote{La volanta o copete es el texto que se ubica sobre el título
de un artículo o nota periodística.} y algunos fragmentos del cuerpo de la
noticia. Luego analice los SSNN subrayados, según se muestra en el ejemplo
resuelto.

\begin{consignas}
  \item Realice el encorchetamiento indicando el núcleo y las funciones de los
        restantes constituyentes.
  \item Indique si los modificadores son restrictivos o no.
\end{consignas}

\textbf{Noticia 1}

\begin{quote}\small
\textit{La diva vuelve a Sudamérica}\\
\textbf{(1)~\subrayado{Madonna} ofrecerá un concierto gratuito en
(2)~\subrayado{la playa de Copacabana}}\\
La reina del pop se presentaría en el mes de mayo. Este lunes, los organizadores
del evento anunciarán todos los detalles.

(3)~\subrayado{El columnista Lauro Jardim del diario O Globo} dijo este mes que
el concierto estaba previsto para el 4 de mayo, fecha aún no confirmada por las
autoridades. (4)~\subrayado{Madonna, de 65 años}, inició su gira
\textit{Celebration} el pasado mes de octubre.
\end{quote}

\textbf{Noticia 2}

\begin{quote}\small
\textit{Es (1)~\subrayado{la primera teatrista latinoamericana} y la segunda
mujer en obtenerlo}\\
\textbf{(2)~\subrayado{Lola Arias, ganadora del Premio Internacional Ibsen
2024}}\\
El jurado destacó la «dimensión ética» y (3)~\subrayado{el carácter
«transnacional» de la obra de la artista multidisciplinaria}.
(\url{www.pagina12.com.ar})
\end{quote}

\begin{resuelto}
Noticia 1. (1) a) [\etq{SN} Madonna\etq{N/n}]. b) El núcleo no presenta
modificadores ni complementos. (2) a) [\etq{SN} la\etq{D/Esp} playa\etq{N/n}
[\etq{SP/M} de\etq{P/n} [\etq{SN/Comp} Copacabana]]]. b) \forma{De Copacabana}
es un modificador restrictivo.
\end{resuelto}
\end{ejercicio}

\begin{ejercicio}\label{ej:05-lobos}
Lea atentamente el siguiente texto y resuelva las consignas que se presentan a
continuación, según se muestra en el ejemplo resuelto:

\begin{consignas}
  \item analice la estructura interna de los SSNN subrayados;
  \item determine en qué casos sería posible o no, dentro de este texto,
        alterar el orden de los constituyentes del sintagma;
  \item fundamente su respuesta teniendo en cuenta los efectos de esos posibles
        cambios (si se produce una modificación del significado o si el
        resultado es agramatical).
\end{consignas}

% Texto de divulgación protegido: se incorpora a mano desde el original.
\textoaincorporar{«Lobos. Dato real», adaptado de A.~Litvin (\textit{El libro de
los villanos}, V\&R, 2015), con los SSNN subrayados: \textit{el lobo, un
mamífero carnívoro; varias subespecies; el famoso aullido; los miembros de su
grupo; muy raras ocasiones}}

\begin{resuelto}
a) [\etq{SN} el\etq{D/Esp} lobo\etq{N/n} [\etq{SN/ApE} un\etq{D/Esp}
mamífero\etq{N/n} carnívoro\etq{SA/Mod}]]. b) Dado que el SN \forma{un mamífero
carnívoro} es una aposición explicativa, puede permutarse por el núcleo al que
modifica ---\forma{un mamífero carnívoro, el lobo}--- sin alterar ni la
gramaticalidad ni el sentido de la estructura que integra.
\end{resuelto}
\end{ejercicio}

\subsection{El sintagma adjetival (SA) y su estructura interna}\label{sec:05-sa}

La estructura del SA puede representarse de la siguiente manera:

\begin{exe}
  \ex\label{ex:05-esquema-sa}
      \begin{forest} gramatica
        [SA
          [Esp\\Q / SAdv]
          [A\prima
            [A\prima
              [A]
              [Comp\\SP]]
            [Mod\\SP]]]
      \end{forest}
\end{exe}

\paragraph{El núcleo.} El núcleo del SA es un adjetivo que puede ser de
distintas subclases: calificativo \exref{ex:05-calif}, relacional
\exref{ex:05-relac} o adverbial \exref{ex:05-adverb}:

\begin{exe}
  \ex\label{ex:05-calif}
    \begin{xlist}
      \ex Fue un accidente \textit{estremecedor}.
      \ex Esta semana se estrenó una \textit{interesante} película de aventuras.
    \end{xlist}
  \ex\label{ex:05-relac}
    \begin{xlist}
      \ex\label{ex:05-relac-a} La reforma \textit{educativa} impulsa un cambio
          de planes.
      \ex\label{ex:05-relac-b} Un escritor \textit{italiano} ganó el premio de
          novela este año.
    \end{xlist}
  \ex\label{ex:05-adverb}
    \begin{xlist}
      \ex Sus \textit{frecuentes} quejas me desesperan.
      \ex El \textit{actual} rey de Inglaterra es Carlos III.
    \end{xlist}
\end{exe}

Los adjetivos \termino{calificativos} conforman el grupo focal o central de la
clase. Son los que expresan propiedades del nombre (\forma{alto, linda, rojo,
amable}). Estos adjetivos admiten graduación (\forma{muy alto, bastante linda,
poco amable}) y suelen tener valores opositivos (\forma{lindo/feo; alto/bajo}).
En cuanto a su posición, pueden ubicarse antes del nombre o después de él
(\forma{un lindo recuerdo / un recuerdo lindo; esa amable mujer / esa mujer
amable}). A diferencia de los calificativos, los adjetivos \termino{relacionales}
no expresan propiedades, sino que restringen el alcance del sustantivo a un
ámbito determinado (\forma{consejo académico, reforma constitucional, educación
cívica}). Como no denotan propiedades, no admiten graduación
(\juicio{*}{\forma{reforma muy constitucional}}), ni manifiestan valores polares,
por lo que no suelen poseer un opuesto. Los relacionales solo pueden aparecer
pospuestos al sustantivo al que modifican (\juicio{*}{\forma{constitucional
reforma}}, \juicio{*}{\forma{cívica educación}}). Como se vio en
\exref{ex:05-relac}, algunos relacionales designan pertenencia a un ámbito
(\forma{político, económico, educativo, sexual}, etc.; v.~\exref{ex:05-relac-a})
o a un lugar determinado (\forma{francés, argentino, bonaerense};
v.~\exref{ex:05-relac-b}); estos últimos son los que se conocen como
\termino{gentilicios}.\pendiente{El original remite a (47) para los
relacionales} Por su parte, los adjetivos \termino{adverbiales} se vinculan con
determinados adverbios referidos a la localización o a la cuantificación
(\forma{antiguo, próximo, habitual, propio, único, simple, último}) y suelen
anteponerse al nombre (\forma{sus frecuentes reclamos / reclama
frecuentemente}). Al igual que los relacionales, no admiten graduación
(\juicio{*}{\forma{su muy / bastante último libro}}).

Los SSAA pueden estar formados por el núcleo solamente (\forma{estaba
\textbf{feliz}, seguía \textbf{triste}}), o bien combinarse con modificadores
(\forma{\textbf{triste} por la noticia}) o complementos (\forma{\textbf{carente}
de recursos, \textbf{apto} para el servicio}) y también pueden recibir
especificadores (\forma{plenamente \textbf{feliz} / muy \textbf{triste}}).

\paragraph{Los complementos.} Como vimos que sucede con los sustantivos, la
mayoría de los adjetivos (\forma{alto, grande, fugaz}, etc.) no necesita
complemento. No obstante, algunos, en función de su significado, exigen un
complemento que se manifiesta mediante un SP encabezado por una preposición
determinada, que puede ser:

\begin{itemize}
  \item \forma{a} (\forma{aficionado, equivalente, propicio, relativo});
  \item \forma{de} (\forma{deseoso, digno, orgulloso, cansado});
  \item \forma{con} (\forma{conforme, contradictorio, satisfecho});
  \item \forma{en} (\forma{experto, abundante, rico});
  \item \forma{para} (\forma{necesario, adecuado, disponible}).
\end{itemize}

Veamos algunos ejemplos:

\begin{exe}
  \ex Mi tío Juan es [muy \textit{aficionado a} la pesca deportiva].
  \ex Todos estaban [\textit{deseosos de} recibir al nuevo integrante de la
      familia].
  \ex El cliente quedó [sumamente \textit{conforme con} la atención].
  \ex Un profesional [\textit{experto en} informática] preparó la encuesta.
  \ex Lea detenidamente el instructivo con los requisitos [\textit{necesarios
      para} el proceso].
\end{exe}

\paragraph{Los modificadores.} Dado que, salvo en los casos que vimos más arriba,
la mayoría de los adjetivos no exige complemento, los SSPP que los acompañan se
consideran modificadores:

\begin{exe}
  \ex Lo vi [feliz \textit{por el encuentro}].
  \ex El vaso está [lleno \textit{hasta la mitad}].
  \ex Estaba [loco \textit{de amor}] por esa mujer.
  \ex Pedrito está [inmenso \textit{de alto}].
\end{exe}

Aunque no es exigido, el SP cuyo complemento es otro adjetivo o un sustantivo
usualmente restringe el alcance del núcleo del SA \parencite{ditullio2014}.

\paragraph{Los especificadores.} Los especificadores del núcleo del SA están
realizados por distintos tipos de cuantificadores, que pueden ser de grado,
\exref{ex:05-muy} y \exref{ex:05-bastante}, o algunos de los terminados en
\forma{-mente} con valor cuantificacional, \exref{ex:05-excesivamente} y
\exref{ex:05-sumamente}.

\begin{exe}
  \ex\label{ex:05-muy} Quedó [\textit{muy} herido por el accidente].
  \ex\label{ex:05-bastante} Es [\textit{bastante} inmaduro para resolver sus
      problemas].
  \ex\label{ex:05-excesivamente} Milena es [\textit{excesivamente} prolija para
      todo], pero su hermano es [\textit{terriblemente} desprolijo].
  \ex\label{ex:05-sumamente} Los inversores estaban [\textit{sumamente}
      interesados en el proyecto].
\end{exe}

También modifican al adjetivo otros adverbios en \forma{-mente}, que pueden
indicar punto de vista \exref{ex:05-pv}, modalidad \exref{ex:05-modalidad} o foco
\exref{ex:05-foco}:

\begin{exe}
  \ex\label{ex:05-pv} Algunos países son [\textit{lingüísticamente /
      étnicamente} heterogéneos].
  \ex\label{ex:05-modalidad} Muchos barrios periféricos son
      [\textit{lamentablemente / indudablemente} peligrosos].
  \ex\label{ex:05-foco} Ese es un riesgo [\textit{específicamente /
      principalmente} económico].
\end{exe}

\begin{ejercicio}\label{ej:05-slogans}
Analice la estructura interna del sintagma de estos viejos eslóganes
publicitarios, según se muestra en el ejemplo resuelto~1.

\begin{exe}
  \ex Demasiado rico para tomarlo despacito. (Nesquik)
  \ex Único en el mundo. / Único. (Fernet Branca)
  \ex Buenos de punta a punta. (Cigarrillos Imparciales)
  \ex Juntos. (Telefé)
\end{exe}

\begin{resuelto}
1. [\etq{SA} demasiado\etq{Q/Esp} [\etq{A\prima} rico\etq{A/n} [\etq{SP/M}
para\etq{P/n} [tomarlo despacito]]]]
\end{resuelto}
\end{ejercicio}

\begin{ejercicio}\label{ej:05-posicion}
Según se muestra en los ejemplos resueltos, en las siguientes series de SSNN
cuyos núcleos son modificados por SSAA:

\begin{consignas}
  \item cambie la posición del modificador;
  \item explique si el cambio produce agramaticalidad, altera el significado o
        si solo produce un efecto estilístico;
  \item en los casos en que cambia el significado, especifique en qué consiste
        la modificación.
\end{consignas}

\begin{exe}
  \ex el funcionario alto
  \ex la crisis económica
  \ex una tarde estupenda
  \ex el discurso político
  \ex ese pobre hombre
  \ex un rico chico
  \ex una enfermedad terrible
  \ex la educación sexual
\end{exe}

\begin{consignas}[resume]
  \item Teniendo en cuenta la resolución de los puntos anteriores, ¿puede
        establecer alguna relación entre lo observado y las subclases del
        adjetivo?
\end{consignas}

\begin{resuelto}[Ejemplos resueltos]
1. a) i) \forma{el funcionario alto} / ii) \forma{el alto funcionario}. b) El
cambio de posición no produce agramaticalidad, pero cambia el significado del A.
c) i) \sig{de elevada estatura} / ii) \sig{importante}. 2. a) i) \forma{la crisis
económica} / ii) \juicio{*}{\forma{la económica crisis}}. b) Si se cambia de
posición, la construcción se vuelve agramatical.
\end{resuelto}
\end{ejercicio}

\begin{ejercicio}\label{ej:05-expansionA}
\begin{consignas}
  \item Realice las expansiones de los siguientes núcleos adjetivos para formar
        sintagmas (considere las respuestas posibles que se muestran abajo para
        el ejemplo~1).
\end{consignas}

\begin{exe}
  \ex \textit{inteligente}
  \ex \textit{propenso}
  \ex \textit{habitual}
  \ex \textit{conforme}
  \ex \textit{deseoso}
\end{exe}

\begin{resuelto}[Respuestas posibles]
1. \forma{sumamente inteligente / muy inteligente para su edad}.
\end{resuelto}

\begin{consignas}[resume]
  \item Teniendo en cuenta la resolución del punto anterior, ¿qué particularidad
        presenta la expansión sintagmática de los adjetivos \forma{propenso,
        conforme} y \forma{deseoso}?
\end{consignas}
\end{ejercicio}

\subsection{El sintagma preposicional (SP) y su estructura interna}\label{sec:05-sp}

La estructura sintagmática del SP puede esquematizarse como sigue:

\begin{exe}
  \ex\label{ex:05-esquema-sp}
      \begin{forest} gramatica
        [SP
          [Esp\\Q]
          [P\prima
            [P]
            [Comp\\SN / SAdv / SP / O]]]
      \end{forest}
\end{exe}

\paragraph{El núcleo.} De todas las clases de palabras que hemos visto, sin duda
la preposición resulta «la clase más conflictiva, en el sentido de que se ha
discutido y se sigue discutiendo su carácter de palabra léxica»
\parencite[56]{giammatteo2009}. Aunque aquí adherimos a este último criterio, no
dejamos de reconocer, como señalan las mencionadas autoras, el carácter atípico
de esta clase de palabras. Para la gramática tradicional, que la preposición no
funcione en forma independiente, es decir, que siempre requiera un complemento,
fue una pauta determinante para considerar que no podía funcionar como núcleo de
un sintagma. Sin embargo, también algunos verbos transitivos y otras palabras,
como los sustantivos eventivos o relacionales, que ya hemos visto en
§\ref{sec:05-sn}, exigen un complemento. Así, por ejemplo, verbos como
\forma{tener} o \forma{regalar}, entre muchos otros, no admiten ir solos:
\juicio{*}{\forma{Yo tengo}}, \juicio{*}{\forma{Juan regaló}}.\footnote{La
omisión del complemento es posible como respuesta a una pregunta, que supone un
contexto previo en el que el complemento ya ha sido mencionado: \forma{¿Quién
trajo bebidas para la reunión? Luis trajo}. Igualmente, algunas preposiciones
permiten la omisión cuando el complemento ya ha sido nombrado: \forma{¿Compraste
té con hierbas? No, sin}.} Por tanto, el hecho de que la preposición siempre
deba ir acompañada por un complemento no impide que sea el núcleo del sintagma
que encabeza. Y esto es así porque, al menos en la mayoría de los casos, las
preposiciones tienen contenido léxico inherente ---\forma{por} (causa),
\forma{para} (fin), \forma{hacia} (dirección), etc.---, por lo que se ubican
dentro de las clases léxicas y funcionan como núcleos de SSPP. No obstante,
también debemos señalar que, en algunos usos que veremos más
adelante,\footnote{El funcionamiento gramatical de la preposición como marca de
función se trata en el \capref{7}.} las preposiciones tienen un funcionamiento
solamente gramatical y poseen escaso significado. En estos casos se considera
que son meras marcas de función que no actúan como núcleos ni determinan la
categoría del sintagma que introducen.

Por último, otra característica que destaca el carácter nuclear de la
preposición es que, al igual que los verbos, restringe semánticamente a su
complemento. Así, por ejemplo, como el complemento de un verbo como
\forma{comer} debe referir a \sig{algo comestible} (\forma{comí pan / postre /
verduras}), una preposición como \forma{hacia} exige que su complemento refiera
a lugar o tiempo ---\forma{hacia el río / el oeste / la noche}---.

Establecido el carácter léxico de la preposición, al menos para la mayoría de
ellas, podemos añadir que son palabras invariables y que, por otra parte, la
característica que las acerca a las palabras funcionales es su carácter
eminentemente relacional. Los SSPP sirven para establecer vínculos de
dependencia con distintas clases de palabras a las que modifican:

\begin{itemize}
  \item sustantivos: \forma{mesa de madera / planta sin flores / estante con
        libros};
  \item adjetivos: \forma{capaz de todo / bueno para nada / sobresaliente en
        matemáticas};
  \item adverbios: \forma{lejos de casa / antes de la madrugada / siempre en
        invierno};
  \item verbos: \forma{estudia en Madrid / piensa en su familia / depende de tu
        decisión / llegó con amigos};\footnote{En el \capref{7} se plantea que
        algunos verbos, como \forma{pensar} y \forma{depender}, exigen
        complementos encabezados por determinadas preposiciones, que se conocen
        como complementos de régimen (CR).}
  \item preposiciones: \forma{por sobre el mar / desde entre los matorrales}.
\end{itemize}

\paragraph{El complemento.} Según ya hemos dicho, dado que las preposiciones
poseen contenido léxico propio, el complemento que seleccionan debe estar acorde
con su significado. Además, la preposición también impone restricciones
formales a su complemento: este no solo es obligatorio, sino que ocupa una
posición fija dentro del sintagma, siempre inmediatamente después del núcleo:
\forma{de colores} / \juicio{*}{\forma{colores de}}.\footnote{Para las
particularidades de la acentuación de la mayor parte de las preposiciones puede
consultarse el \capref{3}, §\ref{sec:03-suprasegmentales}.} Asimismo, al igual
que el verbo exige caso\footnote{La categoría de caso se presentó en el
\capref{4}, §\ref{sec:04-categorias}.} acusativo a su COD, la preposición impone
caso oblicuo o terminal a su complemento cuando este está representado por un
pronombre (\forma{de mí} / \juicio{*}{\forma{de yo}}).

\paragraph{Los especificadores.} Además del complemento, opcionalmente los SSPP
admiten cuantificadores en la posición de especificador:

\begin{exe}
  \ex
    \begin{xlist}
      \ex Poné el cuadro [\textit{más} hacia la derecha].
      \ex {}[\textit{Muy} desde el fondo de mi corazón] debo decir gracias.
    \end{xlist}
\end{exe}

\begin{ejercicio}\label{ej:05-sp-modifica}
Identifique el núcleo al que modifica cada SP y señale a qué clase de palabra
pertenece, según se muestra en el ejemplo resuelto~1.

\begin{exe}
  \ex Plantas sin flores.
  \ex Lo eligieron por sobre los demás.
  \ex Cerca de la oficina.
  \ex Salir con viejos amigos.
  \ex Feliz por sus logros académicos.
\end{exe}

\begin{resuelto}
1. El SP encabezado por \forma{sin} modifica a \forma{plantas}, que es un N.
\end{resuelto}
\end{ejercicio}

\begin{ejercicio}\label{ej:05-de-sp}
Observe los SSPP subrayados de las siguientes oraciones y resuelva las consignas
de abajo, según se muestra en el ejemplo resuelto~1.

\begin{consignas}
  \item Haga la descripción estructural mediante encorchetamiento.
  \item Determine, en función de la palabra a la que modifica el SP, el
        funcionamiento y significado que adopta en cada caso la preposición
        \forma{de}.
\end{consignas}

\begin{exe}
  \ex La casa \subrayado{de mi abuela} es antigua.
  \ex El libro \subrayado{de gramática} resultó útil.
  \ex Me encantan los caramelos \subrayado{de menta}.
  \ex No supieron \subrayado{de aquella mujer} por varios días.
  \ex Los niños, ávidos \subrayado{de conocimiento}, concurrieron al instituto.
\end{exe}

\begin{resuelto}
1. a) [\etq{SP} de\etq{P/n} [\etq{SN/Comp} mi\etq{D/Esp} abuela\etq{N/n}]]. b)
\forma{De} modifica al N \forma{casa} e indica \sig{posesión}.\pendiente{El
original dice «cartera»}
\end{resuelto}
\end{ejercicio}

\begin{ejercicio}\label{ej:05-preposiciones}
Determine los significados de las preposiciones de los SSPP de la lista de
abajo. Use paráfrasis para completar la columna de la derecha, según se muestra
en el ejemplo resuelto.

\begin{center}
\begin{tabular}{@{}r l p{45mm}@{}}
  \toprule
  & SSPP & Significado de la preposición \\
  \midrule
  1.  & asombrados por la noticia & \respuesta{\sig{a causa de}} \\
  2.  & cien metros sobre el nivel del mar & \\
  3.  & el sapo entre los matorrales & \\
  4.  & los problemas sin solución & \\
  5.  & remedios caseros para el empacho & \\
  6.  & el camino hasta su casa & \\
  7.  & tu opinión sobre las instituciones & \\
  8.  & el edicto contra la guerra & \\
  9.  & un viaje hacia bellos lugares & \\
  10. & un pago mediante transferencia bancaria & \\
  11. & la verdad tras su reacción & \\
  \bottomrule
\end{tabular}
\end{center}
\end{ejercicio}

\subsection{El sintagma adverbial (SAdv) y su estructura interna}\label{sec:05-sadv}

De acuerdo con la teoría X\prima, podemos representar la estructura del SAdv
mediante un esquema arbóreo como el siguiente:

\begin{exe}
  \ex\label{ex:05-esquema-sadv}
      \begin{forest} gramatica
        [SAdv
          [Esp\\Q / SN]
          [Adv\prima
            [Adv\prima
              [Adv]
              [Comp\\SP]]
            [Mod\\SN / SP / O]]]
      \end{forest}
\end{exe}

\paragraph{El núcleo.} A diferencia de otras clases léxicas, como el
sustantivo, el adjetivo o el verbo, que manifiestan categorías morfológicas, los
adverbios, al igual que las preposiciones, son palabras invariables. Esto
implica que, cuando se proyectan en estructuras más complejas y funcionan como
núcleo del SAdv, no manifiestan concordancias. Por otra parte, mientras la
preposición, según hemos visto en §\ref{sec:05-sp}, siempre exige complemento
(\forma{por la ruta, sobre el tejado}), el adverbio, por lo general, no lleva
modificadores obligatorios, aunque, opcionalmente, algunos adverbios pueden
admitirlos (\forma{antes de la época actual, ayer por la tarde}).

Desde el punto de vista sintáctico y semántico, los adverbios se presentan como
una extensa y sumamente heterogénea clase, difícil de sistematizar por la
imposibilidad de determinar características comunes a todos sus miembros. Por
eso aquí presentamos clasificaciones que tienen en cuenta distintos criterios,
ya sean semánticos, categoriales o sintácticos.

Semánticamente se distinguen adverbios de:

\begin{itemize}
  \item lugar: \forma{alrededor, cerca, debajo};
  \item tiempo: \forma{siempre, ya, enseguida};
  \item modo: \forma{bien, peor, lentamente};
  \item afirmación: \forma{sí, cierto, obvio};
  \item negación: \forma{no, jamás, tampoco};
  \item duda:\footnote{\label{fn:05-cuantificadores}No incluimos en esta lista a
        los adverbios de cantidad porque hoy día se considera que palabras como
        \forma{muy, poco, bastante} o \forma{demasiado} integran la clase de los
        cuantificadores \parencite{leonetti2007,bosque2009}.}
        \forma{acaso, quizás, posiblemente}.
\end{itemize}

Desde el punto de vista de su naturaleza categorial, la \textit{NGLE}
\parencite[§\,30.2j]{rae2025} distingue dos subclases: adverbios de base léxica
y adverbios gramaticales.

\begin{itemize}
  \item Los \termino{adverbios de base léxica} incluyen la mayoría de los
        terminados en \forma{-mente} (\forma{rápidamente, estrepitosamente}) y
        otros no derivados, como \forma{bien} o \forma{temprano}. Los adverbios
        en \forma{-mente} provienen de adjetivos calificativos y admiten la
        paráfrasis mediante el SP \forma{de modo / de manera} + el A base:
        \forma{astuto > astutamente / de manera astuta}. Algunos adverbios en
        \forma{-mente} tienen una forma corta, sin el sufijo \forma{-mente}:
        \forma{trabajaron duramente / trabajaron duro; habló claramente / habló
        claro; aprendió rápidamente / aprendió rápido}. Estas variantes cortas
        se conocen también como \termino{adverbios adjetivales}, porque tienen la
        forma del adjetivo del que provienen, aunque se diferencian de ellos por
        la ausencia de flexión en género (\forma{Pedro habló claro / Julia habló
        claro}) (v. \capref{4}, §\ref{sec:04-conversion}). El empleo de estos
        adverbios cortos tiene algunas restricciones, ya que no aparecen con
        cualquier tipo de verbo, sino con algunos (\forma{Se portó genial,
        apuntó alto, jugaron sucio, hiló fino}). Varios de estos adverbios son
        propios del español de América: \forma{breve} (= \sig{pronto}),
        \forma{sabroso}, (\forma{oler}) \forma{lindo / rico},
        (\forma{equivocarse}) \forma{feo / fiero}; mientras que otros pertenecen
        al registro peninsular: (\forma{sentar algo}) \forma{fatal},
        (\forma{saberse algo}) \forma{fenomenal} (Bosque,
        1990).\pendiente{En la bibliografía figura Bosque (1989)}
  \item Los \termino{adverbios gramaticales} forman grupos cerrados. Se
        reconocen distintas subclases: \termino{demostrativos} (\forma{aquí,
        hoy, así}); \termino{identificativos} o \termino{referenciales}
        (\forma{antes / después, encima / debajo}); \termino{relativos}
        (\forma{cuando, como, donde}); \termino{enfáticos}\footnote{La subclase
        de adverbios enfáticos incluye los que la \textit{NGLE}
        \parencite[§\,30.2j]{rae2025} distribuye en dos subclases:
        interrogativos y exclamativos.} (\forma{cuándo, cómo, adónde}) y
        \termino{focales} (\forma{precisamente, también, solo, concretamente,
        incluso}).\footnote{Según lo explicado en la
        nota~\ref{fn:05-cuantificadores}, tampoco incluimos entre los adverbios
        gramaticales a los cuantificacionales.}\pendiente{El original da
        \textit{quien} entre los relativos}
\end{itemize}

Según su funcionamiento sintáctico, \textcite[724]{kovacci1999} distinguió dos
grupos de adverbios:

\begin{itemize}
  \item los \termino{adverbios nucleares}, que son los que modifican a un
        adjetivo \exref{ex:05-nucleares-a}, a otro adverbio
        \exref{ex:05-nucleares-b} o tienen funciones circunstanciales dentro del
        SV\footnote{Las funciones sintácticas dentro del SV se estudian en el
        \capref{7}.} \exref{ex:05-nucleares-c}:
\end{itemize}

\begin{exe}
  \ex\label{ex:05-nucleares}
    \begin{xlist}
      \ex\label{ex:05-nucleares-a} Es \textit{relativamente} alto.
      \ex\label{ex:05-nucleares-b} Vive \textit{extremadamente} lejos.
      \ex\label{ex:05-nucleares-c} Llegaron muy \textit{rápidamente}.
    \end{xlist}
\end{exe}

\begin{itemize}
  \item los \termino{adverbios periféricos}, que modifican a toda la oración o a
        la modalidad\footnote{La modalidad y el funcionamiento de los adverbios
        periféricos se tratan en el \capref{14}.} \exref{ex:05-perifericos}:
\end{itemize}

\begin{exe}
  \ex\label{ex:05-perifericos} \textit{Probablemente / Lamentablemente}, no
      esperaba esa noticia.
\end{exe}

\paragraph{Los complementos.} Pocos adverbios reciben complementos. Los que
derivan de adjetivos que exigen un complemento, por lo general, no lo heredan
(\forma{propenso a enfermarse} / \juicio{*}{\forma{propensamente a enfermarse}},
\forma{digno de admiración} / \juicio{*}{\forma{dignamente de admiración}})
\parencite{giammatteo2023a}. En cambio, los locativos que provienen de la unión
de un nombre y una preposición (\forma{de-bajo / a-bajo, de-lante / a-delante,
a-dentro, a-fuera, en-cima, a-rriba, al-rededor, en-frente, de-trás}) y algunos
temporales (\forma{antes, después}), que por su naturaleza categorial se ubican
entre los adverbios identificativos o referenciales que vimos más arriba, pueden
recibir complementos. La particularidad de estos adverbios, que son también
conocidos como \termino{adverbios preposicionales}, es que, al igual que las
preposiciones y a diferencia de otros adverbios, son transitivos
\parencite{giammatteo2013}:

\begin{exe}
  \ex
    \begin{xlist}
      \ex Unos amigos de mi hermana se mudaron [\textit{enfrente} de nuestra
          casa].
      \ex Llegaron muy puntuales, [\textit{antes} de la cena].
    \end{xlist}
\end{exe}

No obstante, la relación del adverbio con su complemento es más laxa que la que
mantiene una preposición, con la que el complemento no es omisible; con el
adverbio, en cambio, el complemento puede quedar implícito:

\begin{exe}
  \ex
    \begin{xlist}
      \ex El equipo quedó [afuera (del torneo)].
      \ex No sé lo que pasó porque yo llegué [después (del suceso)].
    \end{xlist}
\end{exe}

\paragraph{Los modificadores.} Un subgrupo de adverbios temporales, locativos y
algunos modales aceptan ser modificados por aposiciones nominales y
preposicionales (\forma{hoy, lunes; allá, en la esquina}), así como también por
oraciones subordinadas introducidas por un pronombre relativo\footnote{Las
subordinadas relativas, que son las introducidas por pronombres relativos como
\forma{que, cuando, donde}, etc., se introducen en el \capref{13}.}
(\forma{ayer, cuando nos vimos; allí, donde nos encontramos}).

\paragraph{Los especificadores.} Operan como especificadores de los adverbios
los cuantificadores de grado (\forma{muy lentamente, sumamente mal}), SSNN con
valor de grado o medida, que van siempre antepuestos (\forma{dos horas más tarde
de la reunión, tres metros delante de la marca}), y los focalizadores
(\forma{solo hoy, recién ahora}). Sin embargo, no todas las subclases de
adverbios admiten especificadores; así, por ejemplo, solo los graduables
aceptan cuantificadores de grado, de modo que podemos decir \forma{muy
lentamente}, pero no \juicio{*}{\forma{muy históricamente}}; \forma{bastante
lejos}, pero no \juicio{*}{\forma{bastante siempre}}.

\begin{ejercicio}\label{ej:05-adverbios}
Lea atentamente los siguientes fragmentos y resuelva las consignas a
continuación, según se muestra en el ejemplo resuelto:

\begin{consignas}
  \item identifique los adverbios, indique su estatus categorial y
        clasifíquelos semánticamente;
  \item indique si son núcleos o modificadores dentro de los sintagmas en los
        que aparecen.
\end{consignas}

% Fragmentos literarios protegidos: se incorporan a mano desde el original.
\textoaincorporar{1. fragmento de Saki, adaptado de \textit{Cuentos escogidos}
(Claridad, 2007), con los adverbios \textit{débilmente, definidamente, más bien,
elaboradamente}; 2. fragmento sobre el percebe de \textit{La naturaleza}
(Sopena, 1981), con \textit{sumamente, dentro}; 3. fragmento de
«Instrucciones para subir una escalera», de J.~Cortázar (\textit{Historias de
cronopios y de famas}, 1962), con \textit{de frente, atrás,
particularmente}}

\begin{resuelto}
a) \forma{débilmente}: adverbio de base léxica de modo. b) Modifica al adjetivo
\forma{iluminado} dentro del SA [\forma{débilmente iluminado}].
\end{resuelto}
\end{ejercicio}

\begin{ejercicio}\label{ej:05-p12}
Analice en forma completa mediante encorchetamiento los siguientes SSNN
extraídos de distintas noticias del diario \textit{Página/12} de Buenos Aires,
según se muestra en el ejemplo resuelto~1. Tenga en cuenta que los nombres
eventivos, relacionales y de representación llevan complemento.

\begin{exe}
  \ex El desecho de otros minerales potencialmente comercializables.
  \ex Una constante dolarización de excedentes económicos.
  \ex Una corriente de presión sobre este segmento del mercado de cambio.
  \ex El ministro de Turismo, un destacado empresario.
  \ex Un dólar financiero para el turista extranjero.
  \ex Un punto de inflexión en ese estilo de funcionamiento político de la
      coalición de gobierno.
\end{exe}

\begin{resuelto}
1. [\etq{SN} el\etq{D/Esp} desecho\etq{N/n} [\etq{SP/M} de\etq{P/n}
[\etq{SN/Comp} otros\etq{D/Esp} minerales\etq{N/n} [\etq{SA/M}
potencialmente\etq{Adv/M} comercializables\etq{A/n}]]]]
\end{resuelto}
\end{ejercicio}

\begin{recuadro}[Reflexión final]
En este capítulo vimos que las palabras, al igual que otras unidades de la
lengua, no solo se disponen de manera lineal, sino que se organizan en
estructuras jerárquicas que reflejan dependencias y relaciones internas. La
identificación de constituyentes sintácticos y la distinción entre palabras
léxicas y funcionales y sus posibilidades de proyección nos permitió
caracterizar a los sintagmas como unidades intermedias, así como también
describir y explicar su conformación estructural. Para ello, nos basamos en la
teoría X\prima, que sostiene que todos los sintagmas comparten una
configuración estructural endocéntrica, en la que el núcleo selecciona
complementos y admite adjuntos y especificadores, en función de factores tanto
léxicos como sintácticos. Nos detuvimos particularmente en la caracterización de
los SSNN, SSAA, SSPP y SSADV, lo cual se ampliará en los próximos capítulos con
la presentación del SV. El ámbito de la sintaxis, en conclusión, no solo aporta
criterios para definir, describir, analizar y fundamentar el funcionamiento y
la organización interna de sus unidades, sino que también ofrece un marco
teórico que permite establecer una conexión entre forma y significado, lo cual,
sin duda, enriquece nuestra comprensión del funcionamiento del sistema de la
lengua.
\end{recuadro}

% --- Secciones finales -------------------------------------------------
\lecturasrecomendadas

\solucionario

\begin{solucion}{\ref{ej:05-segmentar}}
2. [[Nuestra tía] [traía regalos constantemente]]. 3. [[Muchos estudiantes]
[prefieren [la virtualidad]]]. 4. [[Las medidas] [preocupan [a [la
población]]]]. 5. [[Los jugadores] [ganaron [la copa] [en 2022]]]. 6. [[Aquellas
profesoras] [incorporaron [en [sus clases]] [herramientas tecnológicas]]].
\end{solucion}

\begin{solucion}{\ref{ej:05-homonimia}}
En este ejercicio los constituyentes se identifican, pero no se analizan
exhaustivamente. a) 2. i) [Llenó [la bolsa] [de papel]]; ii) [Llenó [la bolsa
de papel]]. 3. i) [Encontraron [el juguete] [roto]]; ii) [Encontraron [el
juguete roto]]. b) 2. En i) se puede considerar que [la bolsa] y [de papel] son
dos constituyentes independientes. Se interpreta que \sig{la bolsa se llenó con
papel}. En ii), [de papel] modifica a [la bolsa], por lo que conforman un solo
constituyente. Se interpreta que \sig{la bolsa está hecha de papel}. 3. En i) se
puede considerar que [el juguete] y [roto] son dos constituyentes
independientes. Se interpreta que \sig{el juguete fue encontrado roto, pero
antes no lo estaba}. En ii), [roto] modifica a [el juguete], por lo que
conforman un solo constituyente. Se interpreta que \sig{se encontró el juguete
que ya estaba roto}. c) El efecto que produce la homonimia estructural es el de
ambigüedad, en este caso, de tipo sintáctico.
\end{solucion}

\begin{solucion}{\ref{ej:05-expansion}}
Respuestas posibles. 2. El renombrado museo de Ciencias Naturales de la ciudad
de La Plata recibió los nuevos ejemplares de restos fósiles encontrados por los
paleontólogos. / El museo de Arte Moderno de Buenos Aires, una institución
pública dependiente del Gobierno de la Ciudad de Buenos Aires, recibió los
ejemplares de las producciones artísticas más recientes. 3. Mis primas llegaron
el sábado a la noche muy cansadas por el viaje.\pendiente{El original dice
«cansados»} / Mis primas llegaron inmediatamente. 4. La última temporada de mi
serie preferida trajo sorpresas que no esperaba. / La temporada de invierno
trajo numerosas sorpresas referidas a la moda. 5. Los vehículos sin el
dispositivo de TelePASE circulaban muy lentamente. / Los vehículos antiguos
circulaban más lentamente.
\end{solucion}

\begin{solucion}{\ref{ej:05-sustitucion}}
2. Reducción: Varios ministros firmaron un acuerdo con el embajador.
Pronominalización: Ellos lo firmaron con él.
\end{solucion}

\begin{solucion}{\ref{ej:05-encorchetar}}
2. [[Los egresados] [buscaron [con entusiasmo] [un salón de fiestas]]]. Ellos
buscaron con entusiasmo un salón de fiestas. Los egresados lo buscaron con
entusiasmo. Los egresados buscaron así un salón de fiestas. Ellos lo buscaron
así. 3. [[Nuestros candidatos] [presentaron [sus propuestas] [a los
contribuyentes]]]. Ellos presentaron sus propuestas a los contribuyentes.
Nuestros candidatos las presentaron a los contribuyentes. Nuestros candidatos
les presentaron sus propuestas. Ellos se\footnote{\label{fn:05-se}La forma
\forma{se} delante de \forma{lo / la} es una variante de realización del clítico
\forma{le}, es decir, es un alomorfo (v. \capref{10}).} las presentaron. 4. [[Los
fuertes temporales] [provocaron [inundaciones] [en toda la ciudad]]].
Provocaron inundaciones en toda la ciudad. Aclaración: al tratarse de un sujeto
no animado, usualmente no se reemplaza por la forma pronominal \forma{ellos},
sino que se prefiere dejar el sujeto tácito.\footnote{La información sobre el
sujeto tácito se desarrolla en el \capref{6}.} Los fuertes temporales las
provocaron en toda la ciudad. Los fuertes temporales provocaron inundaciones
allí.
\end{solucion}

\begin{solucion}{\ref{ej:05-preguntas}}
2. ¿Quiénes llegaron ayer por la tarde a la ciudad para visitarnos? ¿Qué
hicieron mis primos? ¿Cuándo llegaron mis primos a la ciudad para visitarnos?
¿Adónde llegaron mis primos ayer por la tarde para visitarnos? ¿Para qué
llegaron mis primos a la ciudad ayer por la tarde? 3. ¿Quién es muy exigente con
los exámenes? ¿Cómo es la nueva profesora? ¿Con qué es muy exigente la nueva
profesora?
\end{solucion}

\begin{solucion}{\ref{ej:05-bioy}}
¿Con quién paseaba el escultor? ¿Por dónde paseaba el escultor con el tirano?
¿Dónde presentó el escultor su última obra? ¿Quién presentó su última obra más
allá del laberinto para los extranjeros ilustres, en el extremo de la alameda de
los filósofos decapitados? ¿Qué presentó el escultor más allá del laberinto para
los extranjeros ilustres, en el extremo de la alameda de los filósofos
decapitados?
\end{solucion}

\begin{solucion}{\ref{ej:05-pares}}
a) 2. i. Eliana encontró [el auto] [de repente] / ii. Eliana encontró [el auto
de su primo]. 3. i. El hombre tocó [al perro] [con la mano] / ii. El hombre tocó
[al perro con sarna]. 4. i. Juan abrió [la puerta] [de una patada] / ii. Juan
abrió [la puerta de la habitación].

b) 2. En i., \forma{de repente} y \forma{el auto} son dos constituyentes porque
se mueven separadamente (\forma{De repente Eliana encontró el auto. Eliana
encontró el auto de repente}); se pronominalizan independientemente
(\forma{Eliana lo encontró de repente. Eliana lo encontró así}); responden a
diferentes preguntas (\forma{¿Qué encontró Eliana de repente? ¿Cómo encontró
Eliana el auto?}). En ii., \forma{el auto de su primo} es un solo constituyente
porque se mueve conjuntamente (\forma{El auto de su primo encontró Eliana});
admite una única pronominalización (\forma{Eliana lo encontró}); responde a una
única pregunta (\forma{¿Qué auto encontró Eliana?}). 3. En i., \forma{al perro} y
\forma{con la mano} son dos constituyentes porque admiten las tres pruebas
propuestas: se mueven separadamente (\forma{El hombre tocó con la mano al perro.
Con la mano el hombre tocó al perro}); se pronominalizan independientemente
(\forma{El hombre lo tocó con la mano. El hombre tocó al perro así}); responden
a diferentes preguntas (\forma{¿A quién tocó el hombre con la mano? ¿Con qué tocó
el hombre al perro?}). En ii., \forma{al perro con sarna} es un solo
constituyente porque se mueve conjuntamente (\forma{Al perro con sarna tocó el
hombre}); admite una sola pronominalización (\forma{El hombre lo tocó}); responde
a una única pregunta (\forma{¿A quién tocó el hombre?}). 4. En i., \forma{la
puerta} y \forma{de una patada} son diferentes constituyentes porque admiten las
tres pruebas propuestas: se mueven separadamente (\forma{De una patada Juan abrió
la puerta. La puerta abrió Juan de una patada}); responden a distintas preguntas
(\forma{¿Qué abrió Juan de una patada? ¿Cómo abrió Juan la puerta?}); se
pronominalizan independientemente (\forma{Juan abrió la puerta así. Juan la
abrió de una patada}). En ii., \forma{la puerta de la habitación} es un
constituyente porque se mueve conjuntamente (\forma{La puerta de la habitación
abrió Juan}), admite una única pronominalización (\forma{Juan la abrió}) y
responde a una única pregunta (\forma{¿Qué abrió Juan?}).
\end{solucion}

\begin{solucion}{\ref{ej:05-nucleo}}
2. a) El núcleo es \forma{bueno}. b) \forma{Bueno} es un A y proyecta un SA. 3.
a) El núcleo es \forma{cerca}. b) \forma{Cerca} es un Adv y proyecta un SAdv. 4.
a) El núcleo es \forma{proporcional}. b) \forma{Proporcional} es un A y proyecta
un SA. 5. a) El núcleo es \forma{para}. b) \forma{Para} es una P y proyecta un
SP. 6. a) El núcleo es \forma{inquietud}. b) \forma{Inquietud} es un N y
proyecta un SN. 7. a) El núcleo es \forma{armario}. b) \forma{Armario} es un N y
proyecta un SN. 8. a) El núcleo es \forma{sin}. b) \forma{Sin} es una P y
proyecta un SP.
\end{solucion}

\begin{solucion}{\ref{ej:05-de}}
2. [\etq{SN} varios\etq{Q/Esp} colaboradores\etq{N/n} expertos\etq{A/M}]. 3.
[\etq{SAdv} arriba\etq{Adv/n}\footnote{\forma{Arriba} pertenece al subgrupo de
adverbios denominados identificativos o referenciales, que son transitivos y
llevan un complemento. Los tratamos en §\ref{sec:05-sadv}.} [\etq{SP/Comp}
de\etq{P/n} [\etq{SN/Comp} la\etq{D/Esp} heladera\etq{N/n}]]]. 4. [\etq{SP}
sin\etq{P/n} [\etq{SN/Comp} comentarios\etq{N/n}]].
\end{solucion}

\begin{solucion}{\ref{ej:05-comp-adj}}
2. a) El núcleo es \forma{feria} y proyecta un SN. b) \forma{Del libro} es un SP
no exigido por el núcleo, que funciona como M. 3. a) El núcleo es
\forma{similar} y proyecta un SA. b) \forma{A los demás casos} es un SP exigido
por el núcleo y funciona como Comp. 4. a) El núcleo es \forma{por} y proyecta un
SP. b) \forma{La izquierda} es un SN exigido por el núcleo y funciona como Comp.
5. a) El núcleo es \forma{bajo} y proyecta un SP. b) \forma{La mesada} es un SN
que funciona como Comp, ya que es exigido por el núcleo. 6. a) El núcleo es
\forma{antes} y proyecta un SAdv. b) \forma{De las diez} es un SP exigido por el
núcleo y funciona como Comp.
\end{solucion}

\begin{solucion}{\ref{ej:05-titulos}}
3. a) Es un SN porque el núcleo es el sustantivo \forma{jardín}. b) El D
\forma{el} es un especificador y \forma{de los cerezos} es un SP no requerido
por el núcleo, un modificador adjunto. 4. a) Es un SN porque el núcleo es el
sustantivo \forma{años}. b) El Q \forma{cien} es un especificador y \forma{de
soledad} es un SP no requerido por el núcleo, un modificador adjunto. 5. a) Es
un SAdv porque el núcleo es el adverbio \forma{ágilmente}. b) El núcleo no
presenta modificadores ni complementos. 6. a) Es un SN porque el núcleo es el
sustantivo \forma{verano}. b) El D \forma{ese} es un especificador y \forma{a
oscuras} es un SP no requerido por el núcleo, un modificador adjunto. 7. a) Es
un SA porque el núcleo es el adjetivo \forma{terrenal}. b) El núcleo no presenta
modificadores ni complementos. 8. a) Es un SP porque el núcleo es la preposición
\forma{desde}. b) \forma{La ventana} es un SN exigido por el núcleo, un Comp.
\end{solucion}

\begin{solucion}{\ref{ej:05-vf}}
2. F(also). Tanto las palabras léxicas como las gramaticales pueden ser núcleos
de sintagmas. 3. F(also). El núcleo mantiene diferentes relaciones con los
constituyentes del sintagma: los complementos son los exigidos por el núcleo,
pero los adjuntos son opcionales y los especificadores determinan o cuantifican
al núcleo y a sus complementos y modificadores. 4. V(erdadero). Los complementos
tienen una relación de dependencia estructural con el núcleo y, además,
completan su significado.
\end{solucion}

\begin{solucion}{\ref{ej:05-buenaformacion}}
2. Aunque los sustantivos propios no necesitan ir precedidos de un
determinativo, un sintagma como \forma{el Mario} puede ser aceptable en ciertas
variedades dialectales para referirse a alguien conocido por todos. 3. Es
gramatical porque el sustantivo común \forma{ingredientes} va precedido de un
determinante que lo especifica y del complemento que requiere por su carácter
relacional. 4. Es gramatical porque los nombres de materia o incontables no
requieren ir precedidos por un determinativo en ciertas posiciones sintácticas,
como la de COD. 5. Los nombres abstractos tienen un comportamiento sintáctico
que se asimila al de los de materia, de modo que pueden ir sin determinativo en
ciertas posiciones sintácticas, como la de complemento
preposicional.\footnote{Para ampliar sobre las diferencias de funcionamiento de
las distintas subclases de sustantivos se puede consultar
\textcite{giammatteo2009}.}
\end{solucion}

\begin{solucion}{\ref{ej:05-ratones}}
{}[\etq{SN} el\etq{D/Esp} queso\etq{N/n}] / [\etq{SN} una\etq{Q} compañía\etq{N/n}
[\etq{SP/Comp} de\etq{P/n} [\etq{SN/Comp} control\etq{N/n} [\etq{SP/Comp}
de\etq{P/n} [\etq{SN/Comp} plagas\etq{N/n}]]]]] / [\etq{SN} los\etq{D/Esp}
alimentos\etq{N/n} [\etq{SP/M} con\etq{P/n} [\etq{SN/Comp} alto\etq{A/M}
contenido\etq{N/n} [\etq{SP/Comp} de\etq{P/n} [\etq{SN/Comp} azúcares]]]]].
Aclaración: \forma{compañía} y \forma{contenido} exigen complemento porque son
nombres relacionales; \forma{control} exige un complemento porque es un nombre
eventivo.
\end{solucion}

\begin{solucion}{\ref{ej:05-titulares}}
a) Noticia 1. (3) [\etq{SN} el\etq{D/Esp} columnista\etq{N/n} [\etq{SN/ApR}
[\etq{N/n} Lauro Jardim]] [\etq{SP/Comp} de\etq{P/n} [\etq{SN/Comp}
el\etq{D/Esp} diario\etq{N/n} [\etq{SN/ApR} O Globo]]]]. \forma{Columnista} exige
complemento porque es un nombre relacional. (4) [\etq{SN} Madonna\etq{N/n},
[\etq{SP/M} de\etq{P/n} [\etq{SN/Comp} 65\etq{Q/Esp} años\etq{N/n}]]].
Noticia 2. (1) [\etq{SN} la\etq{D/Esp} primera\etq{A/M} teatrista\etq{N/n}
latinoamericana\etq{A/M}]. (2) [\etq{SN} [\etq{N/n} Lola Arias], [\etq{SN/ApE}
ganadora\etq{N/n} [\etq{SP/Comp} de\etq{P/n} [\etq{SN/Comp} el\etq{D/Esp}
Premio\etq{N/n} Internacional\etq{A/M} Ibsen\etq{SN/ApR}
2024\etq{SN/ApR}]]]].\footnote{La secuencia \forma{Premio Internacional Ibsen
2024} también podría considerarse un compuesto sintagmático que actúa como un
nombre propio (v. \capref{4}, §\ref{sec:04-composicion}).} (3) [\etq{SN}
el\etq{D/Esp} carácter\etq{N/n} transnacional\etq{SA/M} [\etq{SP/M} de\etq{P/n}
[\etq{SN/Comp} la\etq{D/Esp} obra\etq{N/n} [\etq{SP/M} de\etq{P/n}
[\etq{SN/Comp} la\etq{D/Esp} artista\etq{N/n} multidisciplinaria\etq{SA/M}]]]]].

b) \forma{Lauro Jardim}, \forma{del diario O Globo}: restrictivos; \forma{de 65
años}: explicativo; \forma{primera} y \forma{latinoamericana}: restrictivos;
\forma{ganadora del Premio Internacional Ibsen 2024}: explicativo;
\forma{transnacional}: restrictivo; \forma{multidisciplinaria}: restrictivo.
\end{solucion}

\begin{solucion}{\ref{ej:05-lobos}}
a) 2. [\etq{SN} varias\etq{Q/Esp} subespecies\etq{N/n}]. 3. [\etq{SN}
el\etq{D/Esp} famoso\etq{SA/M} aullido\etq{N/n}]. 4. [\etq{SN} los\etq{D/Esp}
miembros\etq{N/n} [\etq{SP/Comp} de\etq{P/n} [\etq{SN/Comp} su\etq{D/Esp}
grupo\etq{N/n}]]]. 5. [\etq{SN} [\etq{SA/M} muy\etq{Q/Esp} raras\etq{A/n}]
ocasiones\etq{N/n}].

b) 2. El Q \forma{varias} puede posponerse al núcleo \forma{subespecies} sin
alterar la gramaticalidad del sintagma, aunque aporta un ligero cambio
estilístico. 3. El A \forma{famoso} puede posponerse al N sin alterar la
gramaticalidad del sintagma, aunque pasa a tener valor restrictivo. 4. El D
\forma{los} no puede posponerse al núcleo por razones
fonológicas\footnote{V. \capref{3}, §\ref{sec:03-acento}.} y la anteposición del
complemento \forma{de su grupo} al núcleo produce un
hipérbaton.\footnote{Figura retórica en la que se altera la sintaxis habitual de
una oración.} 5. El SA puede posponerse al núcleo (\forma{ocasiones muy raras}),
adoptando valor restrictivo.
\end{solucion}

\begin{solucion}{\ref{ej:05-slogans}}
2. [\etq{SA} único\etq{A/n} [\etq{SP/M} en\etq{P/n} [\etq{SN/Comp} el\etq{D/Esp}
mundo\etq{N/n}]]] / [\etq{SA} único\etq{A/n}]. 3. [\etq{SA} buenos\etq{A/n}
[\etq{SP/M} de\etq{P/n} [\etq{SAdv/Comp} [\etq{LAdv/n} punta a
punta]]]].\footnote{\forma{De punta a punta} funciona como una unidad porque es
una locución adverbial. Las locuciones adverbiales se tratan en el
\capref{11}.}\pendiente{El análisis separa \textit{de} de la locución que la
nota declara unitaria} 4. [\etq{SA} juntos\etq{A/n}].
\end{solucion}

\begin{solucion}{\ref{ej:05-posicion}}
a) 3. i) \forma{una tarde estupenda} / ii) \forma{una estupenda tarde}. 4. i)
\forma{el discurso político} / ii) \juicio{*}{\forma{el político discurso}}. 5.
i) \forma{ese pobre hombre} / ii) \forma{ese hombre pobre}. 6. i) \forma{un rico
chico} / ii) \forma{un chico rico}. 7. i) \forma{una enfermedad terrible} / ii)
\forma{una terrible enfermedad}. 8. i) \forma{la educación sexual} / ii)
\juicio{*}{\forma{la sexual educación}}.

b) En 2, 4 y 8, el cambio de posición produce agramaticalidad; en 3 y 7, el
cambio de posición no produce ni agramaticalidad ni cambio de significado del A,
solo aporta un matiz estilístico; en 5 y 6, el cambio de posición no produce
agramaticalidad, pero cambia el significado del A.

c) 5. i) \sig{sin importancia, insignificante} / ii) \sig{sin
dinero}.\pendiente{El original invierte i y ii} 6. i) \sig{lindo} / ii)
\sig{adinerado}.

d) Los A calificativos pueden anteponerse o posponerse al núcleo al que
modifican. Por este motivo, en ninguna de las series que los contienen los
sintagmas se convierten en agramaticales. En 5 y 6 el cambio de posición altera
el significado del adjetivo, pero en 3 y 7 no lo altera y solo manifiesta un
matiz estilístico. Los A relacionales no admiten cambio de posición y siempre
deben posponerse al nombre.
\end{solucion}

\begin{solucion}{\ref{ej:05-expansionA}}
a) Respuestas posibles: 2. \forma{propenso a enfermarse / propenso a decir
mentiras}. 3. \forma{muy habitual por las mañanas / sumamente habitual}. 4.
\forma{conforme con su calificación / conforme a la ley}. 5. \forma{deseoso de
verte / deseoso de comer algo dulce}. b) Los adjetivos \forma{propenso,
conforme} y \forma{deseoso} exigen un complemento encabezado por una
preposición: respectivamente, \forma{a}; \forma{con} o \forma{a}; y
\forma{de}.\pendiente{El original dice «con o por» para
\textit{conforme}}
\end{solucion}

\begin{solucion}{\ref{ej:05-sp-modifica}}
2. El SP encabezado por \forma{por} modifica a \forma{eligieron}, que es un
verbo; el SP encabezado por \forma{sobre} es complemento de \forma{por}, que es
una preposición. 3. El SP encabezado por \forma{de} modifica a \forma{cerca},
que es un Adv. 4. El SP encabezado por \forma{con} modifica a \forma{salir}, que
es un V. 5. El SP encabezado por \forma{por} modifica a \forma{feliz}, que es un
A.
\end{solucion}

\begin{solucion}{\ref{ej:05-de-sp}}
a) 2. [\etq{SP} de\etq{P/n} [\etq{SN/Comp} gramática\etq{N/n}]]. 3. [\etq{SP}
de\etq{P/n} [\etq{SN/Comp} menta\etq{N/n}]]. 4. No supieron [\etq{SP/Comp}
de\etq{P/n} [\etq{SN/Comp} aquella\etq{D/Esp} mujer\etq{N/n}]] por varios días.
5. [\etq{SP} de\etq{P/n} [\etq{SN/Comp} conocimiento]].

b) 2. \forma{De} modifica al N \forma{libro} e indica \sig{tema}. 3. \forma{De}
modifica al N \forma{caramelos} e indica \sig{materia o sustancia}. 4.
\forma{De} es una preposición exigida por el V \forma{saber}. La P introduce un
complemento que modifica el significado del V: \forma{saber de (alguien)} es
\sig{tener información sobre esa persona}. 5. \forma{De} es una P exigida por el
adjetivo \forma{ávidos}, que introduce su complemento.
\end{solucion}

\begin{solucion}{\ref{ej:05-preposiciones}}
Respuestas posibles: 2. \sig{por encima de}. 3. \sig{en medio de}. 4.
\sig{carentes de}. 5. \sig{con el fin de}. 6. \sig{en dirección a}. 7.
\sig{acerca de}. 8. \sig{en oposición a}. 9. \sig{con destino a}. 10. \sig{con,
a través de}. 11. \sig{luego de}.
\end{solucion}

\begin{solucion}{\ref{ej:05-adverbios}}
1. \forma{definidamente}: a) adverbio de base léxica de modo; b) es núcleo del
SAdv [\forma{definidamente}], que modifica al V dentro del SV. \forma{más bien}:
a) locución adverbial de modo; b) especificador del Adv \forma{elaboradamente}
en el SAdv [\forma{más bien elaboradamente}]. \forma{elaboradamente}: a)
adverbio de base léxica de modo; b) modifica al A \forma{rígida} en el SA
[\forma{más bien elaboradamente rígida}]. 2. \forma{sumamente}: a) adverbio de
base léxica de modo; b) es especificador del A en el SA [\forma{sumamente
interesante}]. \forma{dentro}: a) adverbio gramatical referencial de lugar; b)
es núcleo del SAdv [\forma{dentro del caparazón del animal adulto}]. 3.
\forma{frente}: a) adverbio gramatical referencial de lugar; b) complemento de
la P \forma{de} en el SP [\forma{de frente}]. \forma{atrás}: a) adverbio
gramatical referencial de lugar; b) complemento de la P \forma{hacia} en el SP
[\forma{hacia atrás}]. \forma{particularmente}: a) adverbio de base léxica de
modo; b) es especificador del A en el SA [\forma{particularmente incómodas}].
\end{solucion}

\begin{solucion}{\ref{ej:05-p12}}
2. [\etq{SN} una\etq{Q/Esp} constante\etq{A/M} dolarización\etq{N/n} [\etq{SP/M}
de\etq{P/n} [\etq{SN/Comp} excedentes\etq{N/n} económicos\etq{A/M}]]].
3. [\etq{SN} una\etq{Q/Esp} corriente\etq{N/n} [\etq{SP/Comp} de\etq{P/n}
[\etq{SN/Comp} presión\etq{N/n} [\etq{SP/Comp} sobre\etq{P/n} [\etq{SN/Comp}
este\etq{D/Esp} segmento\etq{N/n} [\etq{SP/Comp} de\etq{P/n} [\etq{SN/Comp}
el\etq{D/Esp} mercado\etq{N/n} [\etq{SP/Comp} de\etq{P/n} [\etq{SN/Comp}
cambio]]]]]]]]].
4. [\etq{SN} el\etq{D/Esp} ministro\etq{N/n} [\etq{SP/Comp} de\etq{P/n}
[\etq{SN/Comp} Turismo]], [\etq{SN/Ap} un\etq{Q/Esp} destacado\etq{A/M}
empresario\etq{N/n}]].
5. [\etq{SN} un\etq{Q/Esp} dólar\etq{N/n} financiero\etq{A/M} [\etq{SP/M}
para\etq{P/n} [\etq{SN/Comp} el\etq{D/Esp} turista\etq{N/n}
extranjero\etq{A/M}]]].
6. [\etq{SN} un\etq{Q/Esp} punto\etq{N/n} [\etq{SP/M} de\etq{P/n} [\etq{SN/Comp}
inflexión\etq{N/n} [\etq{SP/Comp} en\etq{P/n} [\etq{SN/Comp} ese\etq{D/Esp}
estilo\etq{N/n} [\etq{SP/Comp} de\etq{P/n} [\etq{SN/Comp} funcionamiento\etq{N/n}
político\etq{A/M} [\etq{SP/Comp} de\etq{P/n} [\etq{SN/Comp} la\etq{D/Esp}
coalición\etq{N/n} [\etq{SP/M} de\etq{P/n} [\etq{SN/Comp}
gobierno\etq{N/n}]]]]]]]]]]].\pendiente{El original dice «su gobierno»}
\end{solucion}
